\documentclass{article}

% ------------------------------------------------------------------
% NeurIPS 2026 style. Default (anonymous) option while under review.
% ------------------------------------------------------------------
\usepackage{neurips_2026}

\usepackage[utf8]{inputenc}
\usepackage[T1]{fontenc}
\usepackage{hyperref}
\usepackage{url}
\usepackage{booktabs}
\usepackage{amsmath}
\usepackage{amssymb}
\usepackage{amsthm}
\usepackage{mathtools}
\usepackage{microtype}
\usepackage{graphicx}
\usepackage{enumitem}
\usepackage{natbib}

\theoremstyle{plain}
\newtheorem{theorem}{Theorem}
\newtheorem{proposition}{Proposition}
\newtheorem{lemma}{Lemma}
\newtheorem{corollary}{Corollary}

\theoremstyle{definition}
\newtheorem{definition}{Definition}
\newtheorem{assumption}{Assumption}
\newtheorem{prediction}{Prediction}
\newtheorem{example}{Example}

\theoremstyle{remark}
\newtheorem{remark}{Remark}

\newcommand{\E}{\mathbb{E}}
\newcommand{\Prob}{\mathbb{P}}
\newcommand{\Var}{\operatorname{Var}}
\newcommand{\R}{\mathbb{R}}
\newcommand{\1}{\mathbf{1}}
\newcommand{\cD}{\mathcal{D}}
\newcommand{\cH}{\mathcal{H}}
\newcommand{\cL}{\mathcal{L}}
\newcommand{\cU}{\mathcal{U}}
\newcommand{\cF}{\mathcal{F}}
\newcommand{\indep}{\perp\!\!\!\perp}
\newcommand{\good}{\mathrm{good}}
\newcommand{\bad}{\mathrm{bad}}
\newcommand{\Tr}{\operatorname{tr}}
\newcommand{\vecop}{\operatorname{vec}}
\newcommand{\Om}{\Omega}
\newcommand{\wdc}{\mathrm{wd}}
\newcommand{\Adam}{\mathrm{AdamW}}

\title{When Does Trace Supervision Induce Stepwise Computation?\\
A Finite-Batch Account of an Optimization Transition}

\author{%
  Anonymous Author(s)\\
  Affiliation\\
  Address\\
  \texttt{email}
}

\begin{document}

\maketitle

% ==================================================================
\begin{abstract}
Transformers trained on chain-of-thought ``think traces'' often show a sharp
change in accuracy as the fraction $\rho$ of correct traces in the training
mixture crosses a modest value. We study this in a fully specified synthetic
word-index task and begin with a negative result: the trace is conditionally
independent of the answer given the prompt, so it cannot reduce Bayes risk. The
Bayes ceiling and the chance floor both have closed forms, $\bar A^\star \approx
0.893$ and $\E[1/L]\approx0.157$, and the empirical curve runs between them with
no fitted parameters. The transition is therefore about which computation SGD
finds, not about what is learnable.

We then analyse a scalar trace-copy coordinate under minibatch SGD. Because the
mixture indicator is drawn independently per example, the full mixture
variance---between-component term included---carries a factor $1/B$, and the
critical fraction is $\rho_c(B) = \kappa\sqrt{v_+}/(\mu_0\sqrt B) + O(B^{-1})$:
it vanishes with batch size rather than approaching a constant. The drift
constant is fixed by the task's own combinatorics, $\mu_0 \propto \bar A^\star -
\E[1/L]$, and the proportionality constant cancels from the threshold.

A three-seed sweep confirms the qualitative prediction---$\rho_c$ falls by a
factor of $7.7$ over a $16$-fold range of $B$---but not the exponent: we measure
$-0.72\pm0.05$ against a predicted $-1/2$, with every seed steeper than $-1/2$
and a conservative seed-level interval that nonetheless contains it. We then ask
whether two natural objections explain the gap. Under a quasi-stationary
surrogate, AdamW replaces the linear drift by a saturating function of the
signal-to-noise ratio and changes the qualitative picture in several ways, but
leaves the leading exponent at $-1/2$; and in a linear-attention surrogate the
exponent is unchanged by parameter dimension. Neither accounts for the gap. What
remains is total sample exposure, for which we give a susceptibility identity and
a control that isolates it. We also report a non-monotonicity at $\rho=1$ present
in all fifteen runs, and identify a token-budget confound in our own evaluation
protocol that is a leading candidate explanation for it.
\end{abstract}

% ==================================================================
\section{Introduction}
\label{sec:intro}

Chain-of-thought training supervises a model on intermediate reasoning tokens as
well as a final answer. In practice the traces used for such supervision are
imperfect: some follow a valid line of reasoning, others are irrelevant or
misleading, and in distilled pipelines a nontrivial fraction may be corrupted
outright. A natural question is when a model trained on such a mixture comes to
\emph{use} the trace as a computational scaffold rather than ignoring it.

In small synthetic settings, final-answer accuracy does not rise smoothly with
the fraction $\rho$ of correct traces. There is a band of $\rho$ below which
accuracy sits near a low baseline and above which it jumps to a high plateau. In
our measurements accuracy moves from $27\%$ to $82\%$ as $\rho$ goes from $0.20$
to $0.40$ at batch size $128$. This paper asks what that transition is.

Three explanations present themselves, and the paper is organised around ruling
out the first and quantifying the second.

\begin{enumerate}[leftmargin=*,itemsep=2pt]
\item \textbf{Information.} Correct traces might carry information about the
answer that the prompt does not. Section~\ref{sec:bayes} shows that in this task
they provably do not.
\item \textbf{Optimization.} Correct traces create a systematic gradient toward a
particular circuit, competing against minibatch noise; the transition is where
signal overcomes noise. This is the account we develop, and it makes a
quantitative prediction about batch size that we test.
\item \textbf{Representability.} The trace-using circuit might not exist in the
architecture. Appendix~\ref{app:repr} argues this is not the binding constraint
at the depths we use.
\end{enumerate}

\subsection{A negative result that reframes the question}

In the word-index task a model receives a word $W$ and a queried character $C$,
and must recover the index $I^\star$ of the occurrence of $C$ that generated the
query. Proposition~\ref{prop:indep} shows $I^\star \indep R \mid (W,C)$, so no
predictor can use the trace $R$ to reduce Bayes risk. The reason is one sentence:
\emph{the correct trace is a deterministic function of the word, so it cannot
tell you anything about the word you did not already have}, and the corrupted
trace is built from independent randomness.

Consequently the Bayes accuracy is fixed by the multiplicity structure of random
words and has the closed form $A^\star(\ell) = \frac{q}{\ell}[1-(1-1/q)^{\ell}]$,
giving $\bar A^\star \approx 0.893$ for our configuration, against a chance floor
$\E[1/L]\approx0.157$. Both are matched by experiment without fitting. Any
explanation of the transition appealing to the trace ``containing the answer'' or
``making the function easier to represent'' is ruled out here. What remains is
the optimizer.

\subsection{A finite-batch account, and one prediction}

Section~\ref{sec:scalar} isolates a scalar coordinate $u$ controlling how much
the answer computation routes through the trace. Correct traces produce a drift
proportional to $\rho$; both kinds of trace produce noise proportional to
$B^{-1/2}$. The transition is where drift overtakes noise. Because each example's
correctness indicator is drawn independently, the entire mixture variance carries
a $1/B$ factor, and the threshold is
\begin{equation}
\label{eq:intro-threshold}
\rho_c(B,\kappa) \;=\; \frac{\kappa\sqrt{v_+}}{\mu_0\sqrt B} + O(B^{-1}),
\end{equation}
so there is no batch-size-independent critical fraction: the critical window
collapses to $\rho = 0$ at rate $B^{-1/2}$. We regard this scaling law, not a
single number, as the falsifiable content of the theory. The drift constant is
not free: at symmetric initialization $\mu_0 \propto \bar A^\star - \E[1/L]$, and
the unknown proportionality constant cancels from \eqref{eq:intro-threshold}
(Remark~\ref{rem:scale}).

\subsection{What we find, and what we do about it}

The qualitative prediction holds: over three seeds $\rho_c$ falls by a factor of
$7.7$ as $B$ grows $16$-fold. The exponent does not: we measure $-0.72\pm0.05$
where the theory predicts $-1/2$.

Two objections to the theory are immediate, and Sections~\ref{sec:adam}
and~\ref{sec:matrix} address them. The bound is derived for plain SGD while the
runs use AdamW; and it is derived for one coordinate while the model has
$8.1\times10^5$ parameters. Under explicitly stated surrogate assumptions,
neither changes the leading exponent. AdamW does change the picture in ways worth
recording---the drift becomes a bounded, saturating function of the
signal-to-noise ratio; the per-step noise stops shrinking with $B$ while the
drift grows; decoupled weight decay turns a horizon-dependent budget criterion
into a horizon-independent threshold---but the exponent stays at $-1/2$. In a
linear-attention surrogate, parameter dimension enters the persistence half of
the argument and not the threshold.

We do not claim this leaves exactly one explanation standing. It leaves total
sample exposure as the channel we cannot rule out, with a susceptibility identity
attached (Proposition~\ref{prop:susceptibility}) and a control that isolates it
(Prediction~\ref{pred:diversity}); and it leaves open that the discrepancy is
partly an artefact of three seeds, since the conservative seed-level interval on
the exponent contains $-1/2$.

\subsection{Contributions}

\begin{enumerate}[leftmargin=*,itemsep=2pt]
\item A formalization of the noisy think-trace word-index task, a proof that the
trace is Bayes-irrelevant, and closed forms for the ceiling and floor bracketing
the empirical curve (Section~\ref{sec:bayes}).
\item A finite-batch escape-and-persistence theorem for a scalar trace-copy
coordinate, with correct mixture-variance accounting and a Freedman bound whose
deviation scale is matched to the drift slack, yielding
$\rho_c = \Theta(B^{-1/2})$; and a derivation of the drift constant from the
task's combinatorics, with the coordinate scale shown to cancel
(Section~\ref{sec:scalar}).
\item A quasi-stationary analysis of AdamW on the same coordinate, giving a
saturating drift field $\Phi(r)$ and a threshold whose leading exponent is
$-1/2$, together with a fixed-horizon invariance result for scale-invariant
optimizers (Section~\ref{sec:adam}).
\item A lift of the scalar analysis to the parameter manifold of a linear
attention layer, showing the exponent is unchanged by dimension under stated
alignment assumptions (Section~\ref{sec:matrix}).
\item A three-seed batch-size sweep confirming the qualitative prediction and
measuring $-0.72\pm0.05$; a susceptibility identity relating the residual to
sample exposure; and identification of a token-budget confound in our evaluation
protocol bearing on the $\rho=1$ anomaly (Sections~\ref{sec:experiments}
and~\ref{sec:gap}).
\end{enumerate}

Throughout we mark which statements are theorems from stated hypotheses, which
are surrogate approximations, which are adopted from the literature, and which
are measurements. Table~\ref{tab:notation} collects notation.

\begin{table}[t]
\caption{Principal notation. Symbols local to a single section are defined there.}
\label{tab:notation}
\centering
\small
\begin{tabular}{ll}
\toprule
Symbol & Meaning \\
\midrule
$\Sigma, q$ & alphabet and its size ($q = 26$) \\
$L, \ell$ & word length (random / fixed) \\
$W, C, I^\star$ & word, queried character, gold index; the target is $Y = I^\star$ \\
$K(W,C)$ & multiplicity: number of $j$ with $W_j = C$ \\
$R$, $Z$ & the presented trace; indicator that it is correct \\
$\rho$, $B$, $T$ & correct-trace fraction, batch size, number of steps \\
$u$, $c$ & scalar trace-copy coordinate; its gradient scale \eqref{eq:gradform} \\
$\mu_0,\mu_1$ & lower / upper bound on correct-trace drift \\
$\epsilon_\bad,\bar\epsilon$ & corrupted-trace bias and its bound \\
$v_+$, $V$ & per-example gradient variance bound; mixture bound $v_+ +\rho(1-\rho)\Delta^2$ \\
$\kappa$, $r$ & confidence factor; standardized SNR $\rho\mu_0\sqrt{B/V}$ \\
$\bar A^\star$, $A_{\mathrm{chance}}$ & Bayes accuracy $\approx0.893$; chance $\E[1/L]\approx0.157$ \\
$S_T$ & cumulative step size $\sum_{s\le T}\eta_s$ \\
\bottomrule
\end{tabular}
\end{table}

% ==================================================================
\section{Related work}
\label{sec:related}

\paragraph{Induction heads and copying circuits.}
Trace-copying is closely related to induction heads, the attention circuits
completing patterns $A\,B\ldots A\to B$ \citep{olsson2022incontext}.
\citet{nichani2024how} analyse how a two-layer transformer learns causal
structure by gradient descent, including the dynamics of an induction-head-like
circuit; our scalar analysis is in the same spirit but specialised to a mixture
of correct and corrupted supervision. \citet{chen2024unveiling} give provable
training dynamics for induction heads with an explicit feature-learning phase.
\citet{ekbote2025what} prove a depth separation between one- and two-layer
transformers for induction heads on Markov chains, which we adopt rather than
re-derive (Appendix~\ref{app:repr}).

\paragraph{Phase transitions and emergence.}
\citet{lavie2026phase} develop a Bayesian theory of copy-head emergence with
phase transitions in attention as a function of task statistics. Work on
emergence and sharp capability jumps
\citep{wei2022emergent,schaeffer2023mirage} motivates our insistence on
identifying the reference levels bracketing a transition before interpreting its
sharpness: because we can compute both floor and ceiling in closed form, we can
say what fraction of the available range a transition traverses rather than
reporting an unnormalised jump.

\paragraph{Sample complexity and information exponents.}
The divergence of the required sample size as the drift vanishes is analogous to
results on the information exponent for single- and multi-index models
\citep{benarous2021online,damian2023smoothing,abbe2023sgd}, where escaping a
symmetric initialization requires a sample size scaling with a negative power of
the population correlation. Appendix~\ref{app:samplecomplexity} records the
corresponding statement here.

\paragraph{Learning under label noise.}
Our setting is learning from a mixture of correct and corrupted supervision, but
unlike classical label noise \citep{natarajan2013learning} the corruption
preserves the final answer token and corrupts only auxiliary trace tokens. The
corrupted trace is independent of the answer by construction and therefore has
zero \emph{initial} bias in the relevant gradient statistic, which is what
produces a batch-size-dependent rather than noise-rate-dependent threshold.

\paragraph{Adaptive optimizers and batch-size scaling.}
Section~\ref{sec:adam} analyses AdamW
\citep{kingma2015adam,loshchilov2019decoupled}. The quasi-stationary treatment of
the accumulators is in the spirit of the SDE approximations and square-root
scaling rules of \citet{malladi2022sdes}, and the saturating drift interpolates
between the gradient and sign-descent \citep{bernstein2018signsgd} regimes;
\citet{zhang2020adaptive} give a complementary account of why adaptivity helps
attention models. The notion of a critical batch size beyond which gradient noise
ceases to bind \citep{mccandlish2018empirical} is the closest existing framing,
and our standardized SNR is its analogue for a single circuit under mixture
supervision. The sample-exposure channel we cannot rule out is the regime
described by data-scaling laws \citep{kaplan2020scaling}.

\paragraph{Expressivity of chain-of-thought, and linear attention.}
A separate line studies what intermediate tokens let a fixed-depth transformer
compute \citep{feng2023towards,merrill2024expressive,li2024chain}. Our task is
constructed so this question does not arise: by Proposition~\ref{prop:indep} the
trace carries no information about the answer beyond the prompt.
Section~\ref{sec:matrix} works with the softmax-free attention map
\citep{katharopoulos2020transformers}, the standard vehicle for provable
statements about attention dynamics
\citep{vonoswald2023transformers,zhang2024incontext}; dimensional factors come
from \citet{tropp2011freedman}.

% ==================================================================
\section{Setup}
\label{sec:setup}

Let $\Sigma$ be an alphabet with $|\Sigma| = q$; in experiments $\Sigma$ is the
lowercase Latin alphabet, $q = 26$, and $L\sim\mathrm{Unif}\{4,\dots,10\}$.

\begin{definition}[Word-index task]
\label{def:task}
Draw $L$; draw $W = (W_0,\dots,W_{L-1})$ uniformly from $\Sigma^L$; draw
$I^\star\sim\mathrm{Unif}\{0,\dots,L-1\}$ independently of $W$ given $L$; set
$C := W_{I^\star}$ and the prompt $P := (W;C)$; draw a trace $R$ per
Definition~\ref{def:traces}. The target is $Y := I^\star$.
\end{definition}

The gold index is drawn \emph{first} and the query read off from it, so the
target is not a deterministic function of $(W,C)$ when $C$ occurs more than once.
This is the key design decision. Defining $I^\star$ as $\min\{j : W_j = C\}$
instead would make Bayes accuracy identically $1$ and the analysis of
Section~\ref{sec:bayes} vacuous: there would be no ceiling below $100\%$ against
which to read the empirical plateau.

\begin{example}
\label{ex:instance}
With $W = \texttt{banana}$ and $I^\star = 3$, we have $C = \texttt{a}$, prompt
$\texttt{banana;a}$, and $K = 3$. From the prompt alone, positions $1$, $3$, $5$
are indistinguishable; the best any predictor can do is $1/3$. The correct trace
$\texttt{0b 1a 2n 3a 4n 5a}$ does not help: it is a re-encoding of the word and
still shows three \texttt{a}'s with no marker distinguishing position $3$.
\end{example}

\begin{definition}[Correct and corrupted traces]
\label{def:traces}
The correct trace is $R^{\good}(W) = (0,W_0)(1,W_1)\cdots(L-1,W_{L-1})$. The
corrupted trace is
$R^{\bad} = (\Pi(0),\widetilde W_0)\cdots(\Pi(L-1),\widetilde W_{L-1})$ with
$\Pi$ a uniform random permutation and $\widetilde W_j$ i.i.d.\ uniform on
$\Sigma$, both independent of $W$, of $I^\star$, and of each other.
\end{definition}

The corruption is \emph{shape-preserving}: same length, same segment count, same
alternation of fields, same marginal law of index values. What it lacks is any
relationship between index and character fields, so corruption cannot be detected
from surface statistics.

\begin{definition}[Training mixture]
\label{def:mixture}
For $\rho\in[0,1]$ draw $Z\sim\mathrm{Bernoulli}(\rho)$ \emph{independently for
each example} and set $R = R^{\good}$ if $Z = 1$, else $R = R^{\bad}$. Write
$\cD_\rho$ for the resulting joint law.
\end{definition}

Per-example independence of $Z$ is used essentially in
Section~\ref{sec:variance}: it is what makes the minibatch gradient an average of
$B$ i.i.d.\ mixture draws and hence what places the between-component variance
inside the $1/B$ factor. It is easy to violate in an implementation---by building
homogeneous batches, or shuffling at shard rather than example level---and
Section~\ref{sec:limitations} records it as a scope condition.

\paragraph{Sequence format and loss.}
The training sequence is $W\texttt{;}C\ R\ \texttt{:}\ I^\star\
\texttt{\textbackslash n}$. Under prompt masking, loss is computed on every token
after the prompt, so with $R$ the trace positions and $Y$ the answer positions,
\begin{equation}
\label{eq:lfull}
\cL_{\mathrm{full}}
= \frac{1}{|R| + |Y|}\Big(
\sum_{i \in R}\cL_{\mathrm{CE}}(r_i, r_i^\star)
+ \sum_{j \in Y}\cL_{\mathrm{CE}}(y_j, y_j^\star)\Big).
\end{equation}
Section~\ref{sec:competing} returns to the consequences of supervising trace
tokens.

\paragraph{Two algorithm families.}
Write $\cH_{\mathrm{short}}$ for predictors that compute the answer by a
content-based search over the prompt for positions $j$ with $W_j = C$ together
with a tie-breaking rule, and $\cH_{\mathrm{trace}}$ for predictors that locate a
trace segment whose character field matches $C$ and copy its index field into the
answer logits. By Corollary~\ref{cor:irrelevance} the two share the same
population ceiling whenever the tie-breaking rule is Bayes-optimal. The question
is not which family has lower population risk---they are tied---but which a
finite run with finite batch size discovers, and how that depends on $\rho$.

% ==================================================================
\section{The trace is Bayes-irrelevant, and both reference levels are computable}
\label{sec:bayes}

\begin{proposition}[Conditional independence]
\label{prop:indep}
Under Definitions~\ref{def:task}--\ref{def:mixture}, for every $\rho$,
$I^\star \indep R \mid (W,C)$.
\end{proposition}

\begin{proof}
\emph{Stage 1: $I^\star\indep R\mid W$.} Condition on $W$ and $L$. The correct
trace is a deterministic function of $W$, hence conditionally independent of
everything given $W$. The corrupted trace is a measurable function of
$(\Pi,\widetilde W)$ alone---despite the notation it does not depend on
$W$---and $(\Pi,\widetilde W)\indep I^\star\mid W$. The indicator $Z$ is
independent of everything else. So $R = \varphi(W,\Pi,\widetilde W,Z)$ with all
arguments other than $W$ conditionally independent of $I^\star$ given $W$, giving
$I^\star\indep R\mid W$.

\emph{Stage 2: adding $C$.} Since $C = W_{I^\star}$ is a function of
$(W,I^\star)$, this is not automatic: conditioning on $C$ does change the
conditional law of $I^\star$. By Stage 1, for bounded $f,g$,
$\E[f(R)g(I^\star)\mid W] = \E[f(R)\mid W]\,\E[g(I^\star)\mid W]$. Conditioning
further on $C$ restricts $I^\star$ to $\{j : W_j = C\}$, reweighting its law; but
the conditional law of $R$ given $W$ does not depend on $I^\star$, so
$\Prob(R\in\cdot\mid W,C=c) = \Prob(R\in\cdot\mid W)$ for every $c$. The product
factorisation survives.
\end{proof}

\begin{corollary}[Trace irrelevance]
\label{cor:irrelevance}
$\Prob(I^\star = j\mid W,C,R) = \Prob(I^\star = j\mid W,C)$ for all $j$ and all
$\rho$. The Bayes-optimal predictor from $(W,C,R)$ coincides with that from
$(W,C)$ and achieves identical $0$--$1$ risk.
\end{corollary}

This is why we prove no theorem of the form ``the trace-using predictor has
strictly lower population risk''. Under
Definitions~\ref{def:task}--\ref{def:mixture} such a statement is false.

\paragraph{The reference levels.}
Fix $W,C$ and let $K = |\{j : W_j = C\}|\ge1$. By construction
$\Prob(I^\star = j\mid W,C) = \1\{W_j = C\}/K$, so $I^\star$ is uniform on the
matching positions and the Bayes predictor attains $1/K$---achieved by any rule
selecting a single matching position, including ``pick the first match'', which
is why $\cH_{\mathrm{short}}$ can be Bayes-optimal.

\begin{proposition}[Closed-form Bayes accuracy]
\label{prop:bayesacc}
For a word of length $\ell$ drawn uniformly from $\Sigma^\ell$,
\begin{equation}
\label{eq:bayesacc}
A^\star(\ell) \;=\; \E[1/K] \;=\; \frac{q}{\ell}\Big[1 - \big(1-\tfrac1q\big)^{\ell}\Big].
\end{equation}
\end{proposition}

\begin{proof}
Condition on $I^\star = i$. The other $\ell-1$ characters are i.i.d.\ uniform and
independent of $C = W_i$, so the number of other occurrences is
$X\sim\mathrm{Binomial}(\ell-1,p)$, $p = 1/q$, and $K = 1+X$. Using
$\frac{1}{1+x}\binom{\ell-1}{x} = \frac1\ell\binom{\ell}{x+1}$, which follows by
writing both sides in factorials,
\[
\E\Big[\frac1{1+X}\Big]
= \frac1\ell\sum_{x=0}^{\ell-1}\binom{\ell}{x+1}p^{x}(1-p)^{\ell-1-x}
= \frac{1}{\ell p}\Big[\sum_{k=0}^{\ell}\binom{\ell}{k}p^{k}(1-p)^{\ell-k} - (1-p)^\ell\Big]
= \frac{1-(1-p)^{\ell}}{\ell p}.
\qedhere
\]
\end{proof}

\begin{corollary}[Levels for our configuration]
\label{cor:levels}
For $q = 26$ and $L\sim\mathrm{Unif}\{4,\dots,10\}$,
\begin{equation}
\label{eq:levels}
\bar A^\star = \tfrac17\textstyle\sum_{\ell=4}^{10} A^\star(\ell) \approx 0.893,
\qquad
A_{\mathrm{chance}} = \E[1/L] \approx 0.157 .
\end{equation}
\end{corollary}

$A_{\mathrm{chance}}$ is the accuracy of guessing uniformly among the $L$
positions, available to any model that has learned only the output format. It is
also the accuracy of the best \emph{constant} predictor, since
$\Prob(I^\star = 0) = \E[1/L]$ and position $0$ is admissible at every length.
The equivalence is useful when generalising to other tasks
(Appendix~\ref{app:family}), where the two baselines differ and the latter is the
right one. Per-length values of $A^\star(\ell)$ are $0.944, 0.926, 0.909, 0.892,
0.875, 0.859, 0.844$ for $\ell = 4,\dots,10$.

\begin{remark}[A genuine information gap needs a different construction]
\label{rem:gap}
Since $R^{\good}$ is a function of $W$ alone, no encoding of it beats the
prompt-only ceiling. To obtain a task where trace access strictly improves Bayes
accuracy, $R^{\good}$ must depend on $I^\star$---the trace must contain the
answer, as a genuine reasoning trace does---while corruption destroys that
dependence. Shrinking $q$ makes the gap large: at $q = 4,\ \ell = 10$ the
prompt-only ceiling is $\approx0.377$, leaving $62$ points for the trace to
recover. We regard this as the most important follow-up to this work.
\end{remark}

% ==================================================================
\section{A scalar theory of the finite-batch threshold}
\label{sec:scalar}

\subsection{The coordinate, and a sign convention}

We isolate a scalar parameter $u$ entering the model only through a logit at the
answer position associated with trace-based copying, so that increasing $u$
increases the log-odds assigned to the index encoded in the attended trace
segment. We study $u_{t+1} = u_t - \eta_t\bar g_t$ with
$\bar g_t = B^{-1}\sum_b \hat g_t^{(b)}$, the minibatch drawn fresh from
$\cD_\rho$ with each $Z$ drawn independently. Since the update subtracts the
gradient, we state all drift hypotheses on $h_t := -\bar g_t$, so that positive
drift means trace-copying is reinforced.

\begin{assumption}[Drift, variance and tails]
\label{ass:main}
There are a neighbourhood $\cU = (-a,a)$ of $u_0 = 0$, constants
$0<\mu_0\le\mu_1<\infty$ and $v_+ > 0$, and
$\epsilon_\bad:\cU\to[0,\infty)$ with $\sup_\cU\epsilon_\bad\le\bar\epsilon$,
such that for all $u\in\cU$: (1) $\mu_0\le\E[-\hat g^{\good}(u)]\le\mu_1$;
(2) $|\E[-\hat g^{\bad}(u)]|\le\epsilon_\bad(u)$;
(3) $\Var(\hat g^{\good}),\Var(\hat g^{\bad})\le v_+$;
(4) the centred per-example gradient is sub-exponential,
$\E\exp(\lambda(\hat g-\E\hat g))\le\exp(\lambda^2v_+/2)$ for $|\lambda|\le1/K_e$.
We call $\bar\epsilon = 0$ the \emph{zero-bias regime}.
\end{assumption}

Item (1) is two-sided deliberately: a one-sided lower bound proves escape but
cannot prevent it, so it cannot support a persistence statement. Item (4) is
stated as a primitive so both Theorem~\ref{thm:escape} and
Appendix~\ref{app:samplecomplexity} invoke the same condition.
Sections~\ref{sec:drift}--\ref{sec:bias} compute $\mu_0$, $v_+$ and
$\epsilon_\bad$ for this task rather than leaving them free.

\subsection{Where the \texorpdfstring{$1/B$}{1/B} goes}
\label{sec:variance}

This subsection is short but carries the main structural claim. Let
$\hat g = Z\hat g^{\good} + (1-Z)\hat g^{\bad}$. By the law of total variance,
\begin{equation}
\label{eq:totalvar}
\Var(\hat g\mid u)
= \underbrace{\rho\Var(\hat g^{\good}) + (1-\rho)\Var(\hat g^{\bad})}_{\text{within}}
+ \underbrace{\rho(1-\rho)\big(\E\hat g^{\good}-\E\hat g^{\bad}\big)^2}_{\text{between}} .
\end{equation}
Writing $\Delta := \mu_1 + \bar\epsilon$ for an upper bound on the mean gap---the
\emph{upper} bound $\mu_1$ is correct here because we are bounding a variance
from above---Assumption~\ref{ass:main} gives
$\Var(\hat g\mid u)\le v_+ + \rho(1-\rho)\Delta^2 =: V$.

Now the decisive step. Because Definition~\ref{def:mixture} draws $Z$
independently per example, the $B$ summands in $\bar g_t$ are i.i.d.\ draws from
the mixture law, so
\begin{equation}
\label{eq:batchvar}
\Var(\bar g_t\mid u_t) \;\le\; \frac{V}{B}
\;=\;\frac1B\Big[v_+ + \rho(1-\rho)\Delta^2\Big],
\end{equation}
with the between-component term \emph{inside} the $1/B$.

\paragraph{Why this is the crux.}
Had a single indicator been shared across a whole minibatch, $\bar g_t$ would be
an average of variables i.i.d.\ only \emph{conditionally} on $Z$, and the law of
total variance applied to $\bar g_t$ would read
$\Var(\bar g_t) = B^{-1}[\text{within}] + \rho(1-\rho)\Delta^2$, with no $1/B$ on
the second term. The drift-dominance condition would then give a strictly
positive interior threshold $\kappa^2/(1+\kappa^2)$ that does not vanish with
$B$. The qualitative prediction of this paper turns on this modelling choice.

\subsection{The threshold and its scaling}

In the zero-bias regime the drift is $\E[h_t\mid u_t]\ge\rho\mu_0$. Requiring it
to exceed the noise standard deviation by a factor $\kappa$ gives
$\rho\mu_0 \ge \kappa\sqrt{V/B}$, i.e.\ after squaring and collecting powers,
\begin{equation}
\label{eq:quadratic}
\rho^2\Big(\mu_0^2 + \tfrac{\kappa^2\mu_1^2}{B}\Big)
\;-\; \rho\,\tfrac{\kappa^2\mu_1^2}{B}
\;-\; \tfrac{\kappa^2 v_+}{B} \;\ge\; 0 .
\end{equation}
The coefficient of $\rho^2$ tends to $\mu_0^2$: the mixture term contributes at
order $1/B$, not order $1$, which is what makes the threshold vanish. Let
$\rho_c(B,\kappa)$ be the unique nonnegative root of \eqref{eq:quadratic} with
equality; it is well defined because with
$A := \mu_0^2+\kappa^2\mu_1^2/B$, $b := \kappa^2\mu_1^2/B$,
$c_q := \kappa^2v_+/B$ the product of roots $-c_q/A$ is negative.

\begin{theorem}[Finite-batch threshold scaling]
\label{thm:scaling}
If $\mu_0 > 0$, $v_+ > 0$ and $\mu_1 <\infty$, then
\begin{equation}
\label{eq:scaling}
\rho_c(B,\kappa)
= \frac{\kappa\sqrt{v_+}}{\mu_0\sqrt B}
+ \frac{\kappa^2\mu_1^2}{2\mu_0^2\,B}
+ O(B^{-3/2}),
\end{equation}
so $\rho_c = \Theta(B^{-1/2})$ and $\rho_c\to0$ as $B\to\infty$.
\end{theorem}

\begin{proof}
The nonnegative root is $\rho_c = (b+\sqrt{b^2+4Ac_q})/(2A)$. The discriminant is
$b^2+4Ac_q = 4\mu_0^2\kappa^2v_+B^{-1}(1+O(B^{-1}))$, since the $B^{-2}$ terms
are lower order, so $\sqrt{b^2+4Ac_q} = 2\mu_0\kappa\sqrt{v_+/B}+O(B^{-3/2})$.
Expanding $1/(2A) = (2\mu_0^2)^{-1}(1-\kappa^2\mu_1^2\mu_0^{-2}B^{-1}+O(B^{-2}))$
and multiplying gives leading term $\kappa\sqrt{v_+}/(\mu_0\sqrt B)$, next term
$\kappa^2\mu_1^2/(2\mu_0^2B)$, and cross terms $O(B^{-3/2})$.
\end{proof}

Equation \eqref{eq:scaling} is the falsifiable content of the scalar theory:
there is no batch-size-independent critical trace fraction in the zero-bias
regime.

\subsection{A two-sided finite-time theorem}
\label{sec:escape}

The drift-dominance condition is a per-step criterion; it does not by itself say
the trajectory leaves $\cU$, which additionally requires enough cumulative step
size.

\begin{theorem}[Escape and persistence]
\label{thm:escape}
Let Assumption~\ref{ass:main} hold with $\bar\epsilon = 0$, let
$S_T = \sum_{s\le T}\eta_s$ and $\sum_{s\le T}\eta_s^2\le C_\eta$, and set
$\rho^{+}_{\mathrm{esc}} := a/(\mu_0S_T)$ and
$\rho^{-}_{\mathrm{per}} := a/(\mu_1S_T)$. Starting from $u_0 = 0$:
\begin{enumerate}[leftmargin=*,itemsep=2pt]
\item[(i)] If $\rho > \rho^{+}_{\mathrm{esc}}$ then, with
$\lambda := \rho\mu_0S_T - a > 0$, the trajectory exits $\cU$ in the
trace-reinforcing direction by time $T$ with probability at least
$1 - \exp\big(-\lambda^2/[2(C_\eta V/B + K_e\lambda/3)]\big)$.
\item[(ii)] If $\rho < \rho^{-}_{\mathrm{per}}$ then, with
$\lambda' := a - \rho\mu_1S_T > 0$, the trajectory remains in $\cU$ for all
$t\le T$ with probability at least
$1 - \exp\big(-\lambda'^2/[2(C_\eta V/B + K_e\lambda'/3)]\big)$.
\end{enumerate}
\end{theorem}

The proof is in Appendix~\ref{app:escape}. Two features matter. First, the
deviation scale $\lambda$ in the Freedman bound is set equal to the actual slack
between cumulative drift and escape radius; choosing instead an unrelated scale
such as $\kappa\sqrt{C_\eta V/B}$ yields an exponent constant in both $B$ and
$\rho$, and hence no concentration statement about the transition at all. Second,
the interval $(\rho^{-}_{\mathrm{per}},\rho^{+}_{\mathrm{esc}})$ is not a proved
transition region but an \emph{undecided} one, in which neither sufficient
condition applies; it is nonempty exactly when $\mu_0<\mu_1$.

\begin{remark}[Two thresholds, and when they agree]
\label{rem:twothresholds}
Theorems~\ref{thm:scaling} and~\ref{thm:escape} answer different questions.
$\rho_c(B,\kappa) = \Theta(B^{-1/2})$ is a per-step signal-to-noise criterion,
depending on $B$ and not the horizon. $\rho^{+}_{\mathrm{esc}} = a/(\mu_0S_T)$ is
a budget criterion, depending on the horizon and not on $B$: with enough
cumulative step size any positive $\rho$ eventually escapes. They agree in
scaling only when $S_T\asymp a\sqrt B/(\kappa\sqrt{v_+})$. This is why
Section~\ref{sec:experiments} holds the number of \emph{steps} fixed rather than
epochs, and it is a confound we return to in Section~\ref{sec:gap}.
\end{remark}

\subsection{The drift constant is fixed by the Bayes structure}
\label{sec:drift}

A threshold stated in unspecified constants is not testable, so we compute them.
Consider a readout at the answer position attending to trace segments whose
character field matches $C$, and let $J$ be the index field of the attended
segment. For a logit parameter $u$ boosting the attended index, the
cross-entropy gradient is
\begin{equation}
\label{eq:gradform}
-\hat g(u) \;=\; c\,\big[\1\{J = I^\star\} - p_u(J)\big],
\qquad c > 0,
\end{equation}
the familiar softmax--cross-entropy residual, with $c = \partial\ell_J/\partial u$
the sensitivity of the attended logit to $u$. The interpretation is that the
trace-copy direction is reinforced exactly to the extent that copying the
attended index would have been right more often than the model's current
confidence in it.

The following assumption makes explicit what is being assumed about how attention
selects $J$; it is not a consequence of Definition~\ref{def:traces}.

\begin{assumption}[Segment selection]
\label{ass:selection}
On correct traces, conditional on $(W,C,R)$ the selected segment $J$ is
measurable with respect to model and input randomness that is independent of
$I^\star$. On corrupted traces, conditional on $(W,C)$ the selected index field
$J$ is uniform on $\{0,\dots,L-1\}$ and independent of $I^\star$.
\end{assumption}

The first clause says the readout cannot break the $K$-fold tie using information
unavailable to a Bayes predictor, which is exactly the content of
Corollary~\ref{cor:irrelevance} transported to the selection mechanism. The
second is a genuine restriction: a trained model could in principle select
segments using the corrupted characters or positional structure, and we assume it
does not. Both clauses are most plausible near initialization and least plausible
late in training, which bounds the scope of what follows.

\begin{proposition}[Bayes-determined initial drift]
\label{prop:mu0}
Under Assumption~\ref{ass:selection}, on correct traces the segments matching $C$
are exactly the $K(W,C)$ positions with $W_j = C$, so
$\Prob(J = I^\star\mid W,C) = 1/K$ and
\begin{equation}
\label{eq:gooddrift}
\E[-\hat g^{\good}(u)] = c\big(\bar A^\star - \E[p_u(J)]\big),
\qquad
\E[-\hat g^{\bad}(u)] = c\big(\E[1/L] - \E[p_u(J)]\big).
\end{equation}
At a symmetric initialization $u_0$ where $p_{u_0}(j) = 1/L$ for all $j$, the
corrupted-trace drift vanishes and
\begin{equation}
\label{eq:mu0}
\mu_0 \;=\; c\,\big(\bar A^\star - \E[1/L]\big) \;\approx\; 0.736\,c .
\end{equation}
\end{proposition}

\begin{proof}
Immediate from \eqref{eq:gradform}, Assumption~\ref{ass:selection} and
Proposition~\ref{prop:bayesacc}, substituting $\E[p_{u_0}(J)] = \E[1/L]$.
\end{proof}

The initial trace-copy signal is therefore proportional to the gap between the
two reference levels that bracket the empirical curve. Note also that the drift
on correct traces is \emph{capped by the Bayes accuracy}: a correct trace does
not hand the model the answer, it hands the model the same $1/K$ ambiguity the
prompt already has. This is the optimization counterpart of
Corollary~\ref{cor:irrelevance}.

\begin{remark}[The coordinate scale cancels]
\label{rem:scale}
The constant $c$ in \eqref{eq:gradform} depends on how $u$ is normalised and is
not determined by the task, so $\mu_0$ alone is not parameter-free. But
$\hat g$ scales linearly in $c$, hence $v_+$ scales as $c^2$, and the threshold
\eqref{eq:scaling} depends on the two only through $\sqrt{v_+}/\mu_0$, in which
$c$ cancels. The numerical prediction of Section~\ref{sec:experiments} is
therefore unaffected by the choice of coordinate scale. The same is not true of
the budget criterion $\rho^{+}_{\mathrm{esc}} = a/(\mu_0S_T)$, because SGD is not
invariant under reparameterisation of $u$; under AdamW, which is scale-invariant
in the gradient, both criteria are invariant.
\end{remark}

\begin{lemma}[Per-example gradient variance]
\label{lem:variance}
Under \eqref{eq:gradform} and Assumption~\ref{ass:selection} at $u_0$, the
per-example gradient variance under $\cD_\rho$ satisfies
\begin{equation}
\label{eq:variance}
\sigma^2(\rho,\ell)/c^2
\;\le\;
\rho\,\bar A^\star(1-\bar A^\star)
+ (1-\rho)\,\tfrac1\ell\big(1-\tfrac1\ell\big)
+ \rho(1-\rho)\Delta^2/c^2 .
\end{equation}
Since the threshold sits at small $\rho$, the relevant value is the $\rho\to0$
limit $v_+ = c^2\,\E[\tfrac1L(1-\tfrac1L)] \approx 0.1296\,c^2$.
\end{lemma}

\begin{proof}
At $u_0$, $p_{u_0}(J) = 1/L$ deterministically, so within each component the
centred gradient is $c(\1\{J=I^\star\}-\pi)$ with $\pi = \Prob(J = I^\star)$, a
scaled centred Bernoulli of variance $c^2\pi(1-\pi)$. By
Proposition~\ref{prop:mu0}, $\pi = \bar A^\star$ on the good component and
$1/\ell$ on the bad one. Substituting into \eqref{eq:totalvar} and bounding the
between term gives the claim.
\end{proof}

The $1/\ell$ dependence that earlier treatments posit as a hypothesis is here a
consequence of the uniformity of $\Pi$ in Definition~\ref{def:traces}.

\subsection{Bias away from initialization}
\label{sec:bias}

The cancellation in Proposition~\ref{prop:mu0} holds at $u_0$ and need not hold
elsewhere: by \eqref{eq:gooddrift},
\begin{equation}
\label{eq:bias}
\epsilon_\bad(u) \;=\; c\,\big|\E[p_u(J)] - \E[1/L]\big| ,
\end{equation}
which is $0$ at $u_0$ and generally positive away from it. This gives a feedback
loop: correct traces push $u$ up, the model becomes more confident,
$\E[p_u(J)]$ rises above $\E[1/L]$, and corrupted traces begin to exert a
negative pull proportional to $(1-\rho)$.

\begin{proposition}[Interior threshold under nonzero bias]
\label{prop:interior}
If $\bar\epsilon > 0$, taking $B\to\infty$ in
$\rho\mu_0 - (1-\rho)\bar\epsilon \ge \kappa\sqrt{V/B}$ gives
\begin{equation}
\label{eq:interior}
\rho_c^\infty(\bar\epsilon) \;=\; \frac{\bar\epsilon}{\mu_0 + \bar\epsilon} > 0 .
\end{equation}
\end{proposition}

So the threshold is $\Theta(B^{-1/2})$ early in training near $u_0$ and drifts
toward an interior value as $u_t$ moves away. We do not bound how far: doing so
would require an assumption controlling $\E[p_u(J)]$ from above over all of
$\cU$, and a model that becomes overconfident---assigning the attended index more
probability than its accuracy warrants---violates any such bound. Nothing in the
Bayes analysis prevents that, since $\E[p_u(J)]$ is a probability and
$\bar A^\star$ an accuracy. We therefore make no claim of the form
$\rho_c^\infty\le\frac12$.

\begin{remark}[Within the scalar readout, the drift stalls at Bayes]
\label{rem:stall}
By \eqref{eq:gooddrift} the expected trace-copy gradient vanishes when
$\E[p_u(J)] = \bar A^\star$. This is a statement about one coordinate of one
readout, not about stationarity of the full network or about its classification
accuracy. Read narrowly, it is a consistency check: the optimization analysis
does not predict progress above the ceiling the information analysis establishes.
\end{remark}

\begin{remark}[Adversarial corruption]
\label{rem:adversarial}
$\bar\epsilon$ can be made positive by construction, by interpolating between
$R^{\good}$ and $R^{\bad}$ so the corrupted trace points at a plausible but
incorrect index rather than a uniform one. This gives a second experimental axis
orthogonal to $\rho$, along which \eqref{eq:interior} predicts a monotone
interior threshold. We do not implement it.
\end{remark}

% ==================================================================
\section{Does the optimizer change the exponent?}
\label{sec:adam}

Section~\ref{sec:scalar} analyses plain SGD; our runs use AdamW. Since the
headline empirical result is a mismatched exponent, the first question is whether
the optimizer accounts for it. This section develops a quasi-stationary surrogate
for AdamW on the same coordinate. Adaptivity turns out to change the mechanism in
several ways worth recording, and to leave the leading exponent at $-1/2$.

We are explicit that this is a surrogate. The moment calculations below assume
the gradient statistics are locally stationary over a window long compared with
the accumulator timescales, which at $\beta_2 = 0.999$ means $\gg10^3$ steps---a
demanding condition during warmup and across the transition itself, where $u_t$
moves quickly. We use it to ask a comparative question (does adaptivity change
the exponent?), not to derive the training trajectory.

\subsection{AdamW as a three-variable system}

\begin{definition}[AdamW on the trace-copy coordinate]
\label{def:adamw}
With $h_t := -\bar g_t$, AdamW iterates
\begin{align}
m_t &= \beta_1 m_{t-1} + (1-\beta_1)\,h_t, \label{eq:adam-m}\\
\nu_t &= \beta_2 \nu_{t-1} + (1-\beta_2)\,h_t^2, \label{eq:adam-v}\\
u_{t+1} &= u_t + \eta_t\,\frac{\hat m_t}{\sqrt{\hat\nu_t}+\varepsilon_{\mathrm A}}
           \;-\; \eta_t\lambda_\wdc\, u_t ,
\label{eq:adam-u}
\end{align}
with $\hat m_t = m_t/(1-\beta_1^t)$, $\hat\nu_t = \nu_t/(1-\beta_2^t)$,
$m_0 = \nu_0 = 0$, decoupled weight decay $\lambda_\wdc\ge0$, and numerical floor
$\varepsilon_{\mathrm A} > 0$ (written $\varepsilon_{\mathrm A}$ to avoid
collision with the bias $\epsilon_\bad$).
\end{definition}

One feature of \eqref{eq:adam-v} drives what follows: $\nu_t$ tracks the second
moment of the \emph{minibatch} gradient, not the per-example one. It therefore
carries its own factor of $1/B$, and the normalization in \eqref{eq:adam-u}
partially undoes the variance reduction a larger batch buys.

\begin{assumption}[Quasi-stationary regime]
\label{ass:qs}
Fix $u\in\cU$ and a window of length
$T_w\gg\max\{(1-\beta_1)^{-1},(1-\beta_2)^{-1}\}$ over which
$\mu := \E[h_t\mid u]$ and $s^2 := \Var(h_t\mid u)$ are constant and $h_t$ is
independent across steps. Assume in addition a bounded normalized fourth moment,
$\Var(h_t^2)\le C_4(\mu^2+s^2)^2$, and that the accumulator covariance is
negligible at the order retained,
$|\mathrm{Cov}(\hat m_t,\hat\nu_t)| = o\big(\mu\,\Var(\hat\nu_t)^{1/2}\big)$.
\end{assumption}

\begin{lemma}[Accumulator moments]
\label{lem:adammoments}
Under Assumption~\ref{ass:qs}, in the stationary regime
\begin{equation}
\label{eq:adammoments}
\E m_t = \mu,\quad
\Var(m_t) = \tfrac{1-\beta_1}{1+\beta_1}s^2,\quad
\E \nu_t = \mu^2 + s^2,\quad
\Var(\nu_t) = \tfrac{1-\beta_2}{1+\beta_2}\Var(h_t^2),
\end{equation}
and by \eqref{eq:batchvar}, $s^2\le V/B$.
\end{lemma}

\begin{proof}
Unrolling, $m_t = (1-\beta_1)\sum_{k\ge0}\beta_1^kh_{t-k}$, so $\E m_t = \mu$ and
$\Var(m_t) = (1-\beta_1)^2s^2\sum_k\beta_1^{2k} =
\frac{1-\beta_1}{1+\beta_1}s^2$. The same computation with $h_t^2$ gives the
$\nu$ statements. Bias correction removes the $O(\beta_i^t)$ transient.
\end{proof}

\begin{remark}[Momentum is an effective batch multiplier, and only that]
\label{rem:momentum}
The momentum variable has the noise of a batch of size
$B(1+\beta_1)/(1-\beta_1)$: at $\beta_1 = 0.9$, $19B$, a better than four-fold
reduction in gradient-noise standard deviation. This is a large effect on the
constant and none on the exponent, since the factor is $B$-free. The same remark
disposes of heavy-ball momentum as an explanation of the discrepancy.
\end{remark}

\subsection{The drift field}

\begin{proposition}[Saturating drift; quasi-stationary approximation]
\label{prop:adamdrift}
Under Assumption~\ref{ass:qs}, define
\begin{equation}
\label{eq:rdef}
r := \frac{\mu}{s} = \frac{\rho\mu_0\sqrt B}{\sqrt V},
\qquad
e_B := \frac{\varepsilon_{\mathrm A}}{s} = \frac{\varepsilon_{\mathrm A}\sqrt B}{\sqrt V}.
\end{equation}
Then the expected AdamW increment at $u$ satisfies
\begin{equation}
\label{eq:adamphi}
\E\big[u_{t+1}-u_t \mid u_t = u\big]
= \eta_t\Big[\Phi(r,e_B) - \lambda_\wdc u\Big]
+ O\!\Big(\eta_t\sqrt{\tfrac{1-\beta_2}{1+\beta_2}}\Big),
\qquad
\Phi(r,e_B) := \frac{r}{\sqrt{1+r^2}+e_B},
\end{equation}
and the per-step fluctuation has standard deviation
$\eta_t c_1/(\sqrt{1+r^2}+e_B)$ with $c_1 := \sqrt{(1-\beta_1)/(1+\beta_1)}$.
\end{proposition}

\begin{proof}
Write $F(m,\nu) = m/(\sqrt\nu+\varepsilon_{\mathrm A})$ and expand about
$(\mu,\E\hat\nu_t)$. The zeroth-order term is
$\mu/(\sqrt{\mu^2+s^2}+\varepsilon_{\mathrm A})$. The first-order terms vanish on
taking expectations, since $\E\hat m_t = \mu$ and $\E\hat\nu_t$ is the expansion
point. The leading surviving corrections are the covariance term
$\partial^2_{m\nu}F\cdot\mathrm{Cov}(\hat m_t,\hat\nu_t)$, negligible by
Assumption~\ref{ass:qs}, and the curvature term
$\tfrac12\partial^2_{\nu}F\cdot\Var(\hat\nu_t)$. The bounded normalized fourth
moment gives $\Var(\hat\nu_t)/(\E\hat\nu_t)^2 = O((1-\beta_2)/(1+\beta_2))$, so
the relative fluctuation of $\sqrt{\hat\nu_t}$ is
$O(\sqrt{(1-\beta_2)/(1+\beta_2)})$, which is the order we retain. Dividing
numerator and denominator by $s$ gives $\Phi(r,e_B)$; the decay term is exact.
The fluctuation statement follows from $\mathrm{sd}(\hat m_t) = c_1s$ and
treating the denominator as deterministic at the same order.
\end{proof}

We state the error at $O(\sqrt{(1-\beta_2)/(1+\beta_2)})$, matching the
denominator fluctuation, rather than at the square of that rate: showing the
first-order contribution cancels exactly would require controlling the
$\hat m$--$\hat\nu$ covariance, which Assumption~\ref{ass:qs} assumes away rather
than establishes. This is the sense in which the section is a surrogate.

Three consequences are worth isolating.

\paragraph{The drift is nonlinear and bounded.}
Under SGD the drift is $\rho\mu_0$, linear in $\rho$ and unbounded. Under AdamW
it is $\Phi(r,e_B)$, behaving as $r/(1+e_B)$ for $r\ll1$ and saturating at
$1/(1+e_B)$. The optimizer responds to the signal-to-noise ratio of the
trace-copy gradient, not to its magnitude.

\paragraph{The noise stops shrinking with $B$; the drift grows instead.}
Under SGD the per-step noise is $\eta_t\sqrt{V/B}$ and the drift
$\eta_t\rho\mu_0$. Under AdamW the per-step noise is at most $\eta_tc_1$,
independent of $B$, while the drift grows with $B$ through $r\propto\sqrt B$.
Batch size enters the two theories through opposite channels; that the exponent
survives this reversal is the content of Theorem~\ref{thm:adamthreshold}.

\paragraph{Weight decay makes the threshold horizon-independent.}
\begin{proposition}[Weight decay converts a budget criterion into a threshold]
\label{prop:wd}
For $\lambda_\wdc > 0$ the deterministic flow
$\dot u = \Phi(r,e_B)-\lambda_\wdc u$ has the unique attracting fixed point
$u^\star = \Phi/\lambda_\wdc$, so the mean trajectory leaves $\cU$ in the
trace-reinforcing direction if and only if
\begin{equation}
\label{eq:wdcriterion}
\Phi(r,e_B) \;>\; \lambda_\wdc\,a ,
\end{equation}
and otherwise is confined for all time rather than merely to a horizon $T$.
\end{proposition}

This resolves an unsatisfying feature of Theorem~\ref{thm:escape}, which predicts
that lengthening training should eventually wash out any threshold, whereas the
observed curves are stable. Combining the two criteria, the binding condition is
$\Phi(r,e_B)\ge\theta$ with
\begin{equation}
\label{eq:thetadef}
\theta := a\,\max\{\lambda_\wdc,\ S_T^{-1}\},
\end{equation}
the first branch the weight-decay fixed point and the second the displacement
$S_T\Phi\ge a$ needed within the horizon. Both branches are $B$-free.

\begin{theorem}[Critical fraction under AdamW]
\label{thm:adamthreshold}
Under Assumption~\ref{ass:qs} in the zero-bias regime, with $\theta$ as in
\eqref{eq:thetadef} and $\theta(1+e_B) < 1$,
\begin{equation}
\label{eq:adamrhoc}
\rho_c^{\Adam}(B)
= \frac{\sqrt V}{\mu_0}\,\frac{\theta}{\sqrt{1-\theta^2}}\;B^{-1/2}
\;+\; \frac{\theta\,\varepsilon_{\mathrm A}}{\mu_0(1-\theta^2)}
\;+\; O\big(\varepsilon_{\mathrm A}^2\sqrt B\big),
\end{equation}
with local log--log slope
$-\tfrac12\big(1+e_B/\sqrt{1-\theta^2}\big)^{-1} \in (-\tfrac12,0]$.
\end{theorem}

\begin{proof}
Solve $F(r,e_B) := r-\theta\sqrt{1+r^2}-\theta e_B = 0$. At $e_B = 0$,
$r_0 = \theta/\sqrt{1-\theta^2}$. Since $\partial_rF = 1-\theta^2$ at $r_0$ and
$\partial_{e_B}F = -\theta$, the implicit function theorem gives
$dr/de_B = \theta/(1-\theta^2)$, so
$r_c = \theta/\sqrt{1-\theta^2}+\theta e_B/(1-\theta^2)+O(e_B^2)$. Substituting
$\rho_c = r_c\sqrt V/(\mu_0\sqrt B)$ and
$e_B = \varepsilon_{\mathrm A}\sqrt{B/V}$ makes the second term $B$-free.
Writing $\rho_c = c_1'B^{-1/2}+c_0'$ with
$c_0'/(c_1'B^{-1/2}) = e_B/\sqrt{1-\theta^2}$ gives the slope. It is negative and
larger than $-\tfrac12$ because $c_0' > 0$: a $B$-independent additive floor can
only flatten a power law.
\end{proof}

\begin{remark}[The flattening is real but numerically invisible]
\label{rem:epsnumeric}
With $\varepsilon_{\mathrm A} = 10^{-8}$ and $\sqrt{v_+}/c = 0.360$, the
standardized floor at our largest batch is $e_{512}\approx6.3\times10^{-7}$. With
the calibrated $\theta = 0.985$ of Appendix~\ref{app:numbers} the predicted slope
is $-0.5(1-3.7\times10^{-6})$. AdamW's numerical floor bends the law toward
\emph{shallower} decay by a few parts in a million.
\end{remark}

\subsection{An invariance result, and its limits}

Theorem~\ref{thm:adamthreshold} handles AdamW. The following says the conclusion
is not special to AdamW, under hypotheses we state carefully because the natural
informal version is too strong.

\begin{theorem}[Fixed-horizon exponent invariance]
\label{thm:noexp}
Consider $u_{t+1} = u_t + \eta_t\Psi(\bar g_{0:t}; u_{0:t})$ with $\Psi$
measurable and positively homogeneous of degree zero in its gradient arguments
jointly. Assume:
\begin{enumerate}[leftmargin=*,itemsep=1pt,label=(H\arabic*)]
\item the horizon $T$ is fixed as $B$ varies, and minibatches are drawn fresh, so
that conditionally on $\cF_t$ the batch gradient is a mean of $B$ i.i.d.\ draws;
\item the per-example gradient has a third absolute moment bounded uniformly over
$u\in\cU$, and $\Var(\hat g(u))$ is bounded away from $0$ on $\cU$;
\item the threshold is the level set of a criterion functional $\mathcal C$ on
the trajectory law that is $B$-free, continuous in the Gaussian limit, and
\emph{transverse} at the threshold: $\partial_r\mathcal C(r)\ne0$ at $r = r_c$.
\end{enumerate}
Then $\rho_c(B) = r_c\sqrt V\mu_0^{-1}B^{-1/2}(1+O(B^{-1/2}))$, so the asymptotic
log--log slope is $-1/2$.
\end{theorem}

\begin{proof}
Write $\bar g_t = -(\mu+s\zeta_t)$ with $\zeta_t$ standardized. By (H1) and (H2),
conditionally on $\cF_t$ the law of $\zeta_t$ is within Kolmogorov distance
$O(B^{-1/2})$ of $N(0,1)$ by the Berry--Esseen theorem, uniformly over
$u\in\cU$; since $T$ is fixed, composing over steps gives an
$O(TB^{-1/2}) = O(B^{-1/2})$ perturbation of the joint law of
$(\zeta_t)_{t\le T}$. By degree-zero homogeneity applied with $c = 1/s$,
$\Psi(\bar g_{0:t};u_{0:t}) = \Psi(-(r+\zeta_0),\dots,-(r+\zeta_t);u_{0:t})$, so
the trajectory is a fixed measurable functional of $(r,(\zeta_t)_{t\le T})$.
Hence $\mathcal C$ depends on $(\rho,B)$ only through $r$, up to $O(B^{-1/2})$.
By transversality (H3) the level set is $r = r_c + O(B^{-1/2})$; substituting
$r = \rho\mu_0\sqrt B/\sqrt V$ gives the claim.
\end{proof}

Hypothesis (H3) is not cosmetic. Without transversality, an $O(B^{-1/2})$
perturbation of the trajectory law can move a level set by an arbitrary amount,
and the conclusion fails. Hypothesis (H1) is what our fixed-step protocol
supplies and what a fixed-sample-budget protocol would violate---a point we
return to in Section~\ref{sec:gap}. Degree-zero homogeneity covers Adam, AdamW,
RMSProp, signSGD, Lion and normalized SGD; plain SGD is degree one and is covered
by Section~\ref{sec:scalar}, where the two criteria of
Remark~\ref{rem:twothresholds} give slopes $-1/2$ and $0$.

\begin{remark}[Why $-1/2$ is hard to beat through the gradient law]
\label{rem:hardfloor}
For a mean of $B$ i.i.d.\ variables with finite variance,
$\Var(\bar g) = \Var(\hat g)/B$ exactly; no averaging scheme concentrates faster.
Any mechanism operating through the minibatch-gradient law therefore has $\rho$
and $B$ entering through $\rho\sqrt B$, and its thresholds are level sets of that
combination. Two candidate departures have the wrong sign: heavy-tailed
per-example gradients with tail index $\alpha\in(1,2)$ give concentration
$B^{-(1-1/\alpha)}$ with $1-1/\alpha<1/2$, a shallower law; and endogenous bias
does the same by Proposition~\ref{prop:biasflat}. This is a heuristic about the
gradient law, not a theorem about all mechanisms---in particular it says nothing
about channels through which $B$ enters outside that law.
\end{remark}

\begin{proposition}[Bias feedback flattens]
\label{prop:biasflat}
Linearise \eqref{eq:bias} as $\epsilon_\bad(u) = \beta|u|$ with
$\beta := c\,\partial_u\E[p_u(J)]|_{u_0}\ge0$. Under AdamW in the small-$r$
regime, evaluating the net drift at the binding point $u = a$, criterion
\eqref{eq:wdcriterion} reads
$(\rho\mu_0-(1-\rho)\beta a)\sqrt{B/V}\ge\theta$, whence
\begin{equation}
\label{eq:biasflat}
\rho_c(B) = \frac{\beta a + \theta a\sqrt{V/B}}{\mu_0 + \beta a},
\qquad
\frac{d\log\rho_c}{d\log B}\in\Big(-\tfrac12,\ 0\Big].
\end{equation}
\end{proposition}

This settles the sign of the bias mechanism: it raises the threshold at every $B$
and pushes the exponent toward $0$, consistent with
Proposition~\ref{prop:interior}. The same argument applies to the horizon: by
\eqref{eq:thetadef} the horizon enters AdamW only inside the $B$-free constant
$\theta$.

% ==================================================================
\section{Does parameter dimension change the exponent?}
\label{sec:matrix}

The second natural objection is that a one-dimensional analysis cannot speak for
a model with $8.1\times10^5$ parameters. This section lifts the escape and
persistence argument to the parameter manifold of a \emph{linear} attention
layer. The conclusion is that dimension enters the persistence half of the
argument and not the threshold.

We state the scope plainly. This is a surrogate result: it shows dimensionality
alone does not alter the leading exponent in a softmax-free attention layer
satisfying explicit alignment assumptions. It does not remove the
scalar-reduction objection for the four-layer softmax transformer we train.

\begin{definition}[Linear attention block]
\label{def:linattn}
Let $X\in\R^{n\times d}$ be the residual stream and $x_q$ the answer-position
row. A linear attention head computes
$\mathrm{Attn}(X) = W_{OV}X^\top XW_{QK}^\top x_q
= \sum_j\langle W_{QK}^\top x_q, x_j\rangle W_{OV}x_j$, and the answer logits are
$\ell = E^\top\mathrm{Attn}(X)$ for fixed unembedding $E$. Collect the trainable
parameters into $U := (W_{QK},W_{OV})\in\R^D$, $D = 2d^2$, and write
$\cU_a := \{U : \|U-U_0\|_F < a\}$.
\end{definition}

\begin{definition}[Trace-copy direction]
\label{def:tracedir}
With $G^\good(U) := \E[-\nabla_U\cL\mid Z=1]$, set
$\Om^\star := G^\good(U_0)/\|G^\good(U_0)\|_F$ and
$u_t := \langle U_t - U_0,\Om^\star\rangle$.
\end{definition}

\begin{assumption}[Trace-copy alignment]
\label{ass:alignment}
There is $\delta\ge0$ with
$\langle\nabla_U\ell_k,\Om^\star\rangle = c\,\1\{k = J\}+O(\delta)$ on $\cU_a$,
$c>0$: the direction $\Om^\star$ selectively boosts the copied index and is
$\delta$-incoherent with the remaining logit directions.
\end{assumption}

Assumption~\ref{ass:alignment} is an assumption, not a consequence of
Definition~\ref{def:tracedir}: normalising the good-component gradient fixes a
direction but does not by itself make that direction selective for the copied
index.

\begin{proposition}[The scalar gradient form is a projection]
\label{prop:projection}
The per-example matrix gradient is
$-\nabla_U\cL = \sum_k(\1\{k=I^\star\}-p_U(k))\nabla_U\ell_k$, so under
Assumption~\ref{ass:alignment}
\begin{equation}
\label{eq:matproj}
\langle-\nabla_U\cL,\ \Om^\star\rangle
= c\big(\1\{J=I^\star\}-p_U(J)\big) + O(\delta),
\end{equation}
which is \eqref{eq:gradform} up to the $O(\delta)$ remainder. Hence
Proposition~\ref{prop:mu0}, Lemma~\ref{lem:variance} and \eqref{eq:bias} apply to
the projected coordinate, with the same constant $c$.
\end{proposition}

\begin{proof}
The displayed gradient is the softmax--cross-entropy residual in vector form.
Taking the inner product with $\Om^\star$ and applying
Assumption~\ref{ass:alignment} termwise, the $k=J$ term contributes
$c(\1\{J=I^\star\}-p_U(J))$ and the remaining terms contribute $O(\delta)$ each
against coefficients summing in absolute value to at most $2$.
\end{proof}

\begin{assumption}[Matrix drift, covariance and tails]
\label{ass:matrix}
On $\cU_a$: (1) $\E[-\hat G^\good(U)] = \mu(U)\Om^\star+\Xi^\good(U)$ with
$\mu_0\le\mu\le\mu_1$ and $\|\Xi^\good\|_F\le\bar\epsilon_\perp$;
(2) $\E[-\hat G^\bad(U)] = \Xi^\bad(U)$ with $\|\Xi^\bad\|_F\le\epsilon_\bad(U)$;
(3) $\mathrm{Cov}(\vecop\hat G)\preceq\Sigma$, with
$v_\parallel := \langle\Om^\star,\Sigma\Om^\star\rangle\le v_+$ and
$D_{\mathrm{eff}} := \Tr\Sigma/v_+$;
(4) $\langle\hat G,\Om\rangle$ is sub-exponential with parameter $K_e$ for every
unit $\Om$.
\end{assumption}

$D_{\mathrm{eff}}$, not $D$, is what enters: by \eqref{eq:matproj} each
per-example gradient is a sum of a few rank-one features, so $\Sigma$ has rapidly
decaying spectrum.

\begin{theorem}[Matrix escape and persistence]
\label{thm:matrixescape}
Let Assumption~\ref{ass:matrix} hold with
$\epsilon_\bad\equiv0,\bar\epsilon_\perp = 0$, and write
$V_\parallel := v_\parallel+\rho(1-\rho)\mu_1^2$,
$V_\perp := \Tr\Sigma+\rho(1-\rho)\mu_1^2$. Then:
\begin{enumerate}[leftmargin=*,itemsep=2pt]
\item[(i)] if $\rho > a/(\mu_0S_T)$, the trajectory exits $\cU_a$ in the
trace-reinforcing direction by time $T$ with probability at least
$1-\exp(-\lambda^2/[2(C_\eta V_\parallel/B+K_e\lambda/3)])$,
$\lambda := \rho\mu_0S_T-a$; \emph{this bound does not involve $D$};
\item[(ii)] if $\rho < a/(\mu_1S_T)$, the trajectory remains in $\cU_a$ for all
$t\le T$ with probability at least
$1-(D+1)\exp(-\lambda'^2/[2(C_\eta V_\perp/B+K_e\lambda'/3)])$,
$\lambda' := a-\rho\mu_1S_T$.
\end{enumerate}
\end{theorem}

Part (i) is proved by projecting onto $\Om^\star$, after which the process is
exactly the scalar recursion with $V\mapsto V_\parallel$; no union bound over
directions is needed, which is why $D$ is absent. Part (ii) requires controlling
all $D$ coordinates simultaneously and is proved by applying the matrix Freedman
inequality of \citet{tropp2011freedman} to the self-adjoint dilation of the
$\R^D$-valued martingale. Appendix~\ref{app:matrix} gives both arguments,
including the dilation and the variance-process bound; we flag that the
dimensional factor there is the standard one for that inequality and that we have
not optimised it.

\begin{corollary}[Dimension-free exponent]
\label{cor:dimfree}
$\rho_c^{\mathrm{mat}}(B,\kappa) = \kappa\sqrt{v_\parallel}\mu_0^{-1}B^{-1/2}
+O(B^{-1}) = \Theta(B^{-1/2})$, uniformly in $D$. Dimension enters only the
persistence half, inflating the width of the undecided window by
$\sqrt{D_{\mathrm{eff}}}$ and its failure probability by $D+1$.
\end{corollary}

\begin{remark}[Why spreading the signal does not dilute it]
\label{rem:spread}
One might expect distributing the signal over $D$ coordinates to weaken it by
$1/\sqrt D$. It does not, because the noise is projected along with the signal:
the drift along $\Om^\star$ is $\|G^\good\|_F$ regardless of how many coordinates
carry it, while the noise along $\Om^\star$ is a variance in \emph{one}
direction. The ratio governing escape is intrinsically one-dimensional.
\end{remark}

Finally, coordinate-wise adaptivity does not interact with the geometry in a way
that changes the exponent: applying Proposition~\ref{prop:adamdrift}
coordinate-wise and projecting gives drift
$\eta_t\sum_i\Om^\star_i\Phi(r_i,e_{B,i})$, which in the small-$r$ regime is
$\eta_t\sqrt B\rho\mu_0\sum_i\Om^{\star2}_i/\sqrt{\Sigma_{ii}}\,(1+o(1))$, again
a $B$-free level set of $\rho\sqrt B$.

\begin{remark}[One initialization where the law differs]
\label{rem:saddle}
If the readout were degenerate at initialization ($\mu(U_0) = 0$) with drift
linear in displacement along the unstable direction of the bilinear map, escape
would be noise-seeded rather than drift-driven, and
$\rho_c\asymp(\log B+C)/(2\mu_0'S_T)$, \emph{increasing} in $B$: larger batches
supply less of the noise needed to seed a saddle escape. Our task is not in this
regime, since Proposition~\ref{prop:mu0} gives $\mu_0>0$ at the symmetric
initialization.
\end{remark}

% ==================================================================
\section{Experiments}
\label{sec:experiments}

\subsection{Setup}

We train a $4$-layer, $4$-head, width-$128$ causal transformer
($\approx8.1\times10^5$ parameters) with a character-level vocabulary, context
length $64$, AdamW, learning rate $3\times10^{-4}$ with $600$ warmup steps and
cosine decay to $10^{-5}$, weight decay $0.01$, dropout $0.05$, label smoothing
$0.1$, gradient clipping at $1.0$, and $6000$ steps per run. Every $(\rho,B)$
configuration is run at three seeds $(2200,1337,2026)$ controlling initialization
and batch order, for $165$ runs. The training pool contains $1{,}536{,}000$
examples generated once, held fixed, and shared across all runs; for a given
$\rho$ the first $\rho N$ examples use correct traces and the rest corrupted
ones, with correct and corrupted variants of the same instance never both used.

\paragraph{The evaluation metric, stated precisely.}
Evaluation is on $500$ held-out prompts. The model is given the \emph{bare
prompt} $W\texttt{;}C$---the trace is \emph{not} supplied as context---and
generates greedily for up to $35$ new tokens or until a newline. The predicted
answer is then parsed from the generated text by taking the segment after the
final colon and reading the last digit string in it, and is scored by exact
string match against the gold index.

This is a compound, end-to-end metric, not teacher-forced accuracy at the answer
position. It confounds emitting a well-formed trace, reaching the delimiter,
predicting the answer, and generation length. Section~\ref{sec:anomaly} shows the
confound is not hypothetical: at $\ell = 10$ the answer arrives on token $34$ of
a $35$-token budget, so a single spurious token truncates it. We report the
metric we actually computed and flag the teacher-forced measurement as the
primary fix in Section~\ref{sec:limitations}.

\subsection{The transition and its reference levels}

Table~\ref{tab:batchsweep} reports accuracy against $\rho$ at five batch sizes
over three seeds. Reading any column: the low-$\rho$ baseline sits at
$15.3$--$17.0\%$ against a computed $A_{\mathrm{chance}} = 15.7\%$, and the
plateau at $87$--$90\%$ against a computed $\bar A^\star = 89.3\%$. Both levels
are matched with no fitted parameters, and away from the transition the
seed-to-seed standard deviations are below $2.5$ points, so the agreement is not
an artefact of seed selection. Because the metric is end-to-end, attaining the
Bayes ceiling additionally implies that trace generation and parsing cost
essentially nothing in the plateau regime---which makes their apparent cost at
$\rho = 1$ the more striking.

The transition is sharp: at $B = 128$ accuracy moves from $27.2\%$ at
$\rho = 0.20$ to $81.6\%$ at $\rho = 0.40$, from $16\%$ to $90\%$ of the
available range, over a doubling of $\rho$.

\paragraph{Variance peaks inside the transition, as predicted.}
The standard deviations are strongly $\rho$-dependent. At $B = 256$ they run
$1.0, 0.4, 5.8, 18.3, 24.9, 1.9, 1.7,\dots$ as $\rho$ increases, peaking at
$\rho = 0.20$ and collapsing on either side. The $24.9$-point figure comes from
runs at $43.9\%$, $88.1\%$ and $46.2\%$: one seed escaped and two did not. Near
$r\approx1$ the escape event of Theorem~\ref{thm:escape} is genuinely stochastic,
so the outcome should be close to a coin flip rather than a smooth intermediate
value, and it is. We regard this as the most direct evidence in the paper for the
escape-versus-confinement picture, as distinct from a smoothly improving circuit.

\begin{table}[t]
\caption{Accuracy (\%) versus correct-trace fraction $\rho$ at five batch sizes,
steps fixed at $6000$. Mean $\pm$ s.d.\ over three seeds. Chance is $15.7\%$ and
the Bayes ceiling $89.3\%$; neither is fitted. Bold marks cells inside each
column's transition. $\hat\rho_c$ is the logistic fit to the seed-averaged curve
with the $\rho = 1$ row excluded.}
\label{tab:batchsweep}
\centering
\small
\begin{tabular}{lccccc}
\toprule
$\rho$ & $B=32$ & $B=64$ & $B=128$ & $B=256$ & $B=512$ \\
\midrule
0.00 & $16.0_{\pm1.8}$ & $16.3_{\pm2.1}$ & $17.0_{\pm2.3}$ & $16.1_{\pm1.0}$ & $16.4_{\pm0.9}$ \\
0.05 & $16.0_{\pm0.9}$ & $15.3_{\pm1.4}$ & $16.1_{\pm0.2}$ & $16.6_{\pm0.4}$ & $16.7_{\pm1.2}$ \\
0.10 & $16.1_{\pm0.4}$ & $16.4_{\pm0.6}$ & $17.8_{\pm2.0}$ & $20.0_{\pm5.8}$ & $\mathbf{54.3}_{\pm29.6}$ \\
0.15 & $16.3_{\pm1.4}$ & $16.3_{\pm1.4}$ & $21.4_{\pm3.0}$ & $\mathbf{48.4}_{\pm18.3}$ & $\mathbf{85.7}_{\pm7.2}$ \\
0.20 & $16.7_{\pm1.6}$ & $18.0_{\pm1.5}$ & $\mathbf{27.2}_{\pm9.9}$ & $\mathbf{59.4}_{\pm24.9}$ & $88.8_{\pm1.4}$ \\
0.30 & $16.8_{\pm0.2}$ & $20.8_{\pm4.4}$ & $\mathbf{47.9}_{\pm11.6}$ & $88.0_{\pm1.9}$ & $88.0_{\pm2.0}$ \\
0.40 & $19.7_{\pm2.4}$ & $\mathbf{41.2}_{\pm15.0}$ & $\mathbf{81.6}_{\pm11.5}$ & $88.4_{\pm1.7}$ & $90.2_{\pm0.8}$ \\
0.50 & $\mathbf{27.9}_{\pm10.0}$ & $\mathbf{61.7}_{\pm23.5}$ & $87.9_{\pm2.3}$ & $89.1_{\pm1.0}$ & $89.7_{\pm0.8}$ \\
0.60 & $\mathbf{35.7}_{\pm13.1}$ & $\mathbf{73.6}_{\pm8.3}$ & $89.6_{\pm0.1}$ & $89.1_{\pm0.6}$ & $89.3_{\pm0.7}$ \\
0.80 & $\mathbf{56.7}_{\pm11.2}$ & $87.1_{\pm3.3}$ & $87.9_{\pm2.8}$ & $88.1_{\pm2.0}$ & $89.2_{\pm1.0}$ \\
1.00 & $39.5_{\pm14.9}$ & $59.8_{\pm12.7}$ & $61.7_{\pm16.7}$ & $71.7_{\pm3.3}$ & $79.2_{\pm6.5}$ \\
\midrule
$\hat\rho_c$ & 0.759 & 0.465 & 0.305 & 0.175 & 0.099 \\
\bottomrule
\end{tabular}
\end{table}

\subsection{The threshold vanishes with batch size, at a rate we did not predict}
\label{sec:exponent}

We fit the logistic profile
$A(\rho) = A_{\mathrm{chance}} + (\bar A^\star-A_{\mathrm{chance}})/
(1+\exp(-(\rho-\rho_c)/\tau))$ with both levels pinned at their computed values,
leaving only $\rho_c$ and the width $\tau$ free, excluding the $\rho = 1$ row.

\paragraph{The qualitative prediction holds.}
$\hat\rho_c$ falls monotonically from $0.759$ at $B = 32$ to $0.099$ at
$B = 512$---a factor of $7.7$ over a $16$-fold increase---and monotonically in
each seed separately. There is no batch-size-independent critical trace fraction,
and a theory predicting a fixed interior threshold is refuted by these data.

\paragraph{The exponent does not match.}
Regressing $\log\hat\rho_c$ on $\log B$ over all fifteen $(B,\mathrm{seed})$
cells,
\begin{equation}
\label{eq:measured}
\hat\rho_c \propto B^{-0.72\pm0.05},\qquad R^2 = 0.95,
\end{equation}
and the seed-averaged curve gives $-0.73$ with $R^2 = 0.997$, against a predicted
$-1/2$. The midpoint-crossing estimator gives $-0.74$, so this is not an artefact
of the fitting profile. Table~\ref{tab:perseed} gives the per-seed fits, and
Figure~\ref{fig:phase} plots the seed-averaged curves against the two computed
reference levels.

\paragraph{How firmly is $-1/2$ excluded?}
Two analyses of the same three seeds disagree, and the disagreement should be
stated rather than resolved silently. Treating each of the fifteen fitted
thresholds as a point gives $-0.72\pm0.05$ with a $95\%$ interval
$[-0.80,-0.64]$, excluding $-1/2$. Treating the \emph{seed} as the unit of
replication---the conservative choice, since all runs share one data pool---the
three per-seed exponents are $-0.794$, $-0.782$ and $-0.570$, with mean $-0.715$
and standard error $0.073$; with two degrees of freedom the $95\%$ interval is
$[-1.03,-0.40]$, which contains $-1/2$ ($p = 0.10$). We report the first as the
point estimate and the second as the honest uncertainty. What can be said without
qualification is that all three seeds give an exponent steeper than $-1/2$, and
that the third seed's $-0.570$ is driven by one cell: $\hat\rho_c = 0.131$ at
$B = 512$ against $0.089$ for the other two.

\begin{table}[t]
\caption{Per-seed threshold estimates and exponents. Every seed gives a
monotonically decreasing $\hat\rho_c$ and an exponent steeper than $-1/2$; the
spread is nonetheless large enough that a seed-level test does not exclude
$-1/2$.}
\label{tab:perseed}
\centering
\small
\begin{tabular}{lccccccc}
\toprule
Seed & $B{=}32$ & $B{=}64$ & $B{=}128$ & $B{=}256$ & $B{=}512$ & exponent & $R^2$ \\
\midrule
2200 & 0.849 & 0.557 & 0.349 & 0.208 & 0.088 & $-0.794$ & 0.977 \\
1337 & 0.754 & 0.415 & 0.266 & 0.132 & 0.089 & $-0.782$ & 0.994 \\
2026 & 0.649 & 0.420 & 0.301 & 0.198 & 0.131 & $-0.570$ & 0.999 \\
\midrule
mean curve & 0.759 & 0.465 & 0.305 & 0.175 & 0.099 & $-0.730$ & 0.997 \\
\bottomrule
\end{tabular}
\end{table}

\begin{figure}[t]
\centering
\includegraphics[width=0.72\textwidth]{figures/phase_transition_low_variance.png}
\caption{Accuracy versus correct-trace fraction $\rho$. The horizontal reference
levels are the computed chance floor $0.157$ and Bayes ceiling $0.893$ of
Corollary~\ref{cor:levels}; neither is fitted.}
\label{fig:phase}
\end{figure}

\subsection{A non-monotonicity at \texorpdfstring{$\rho=1$}{rho=1}}
\label{sec:anomaly}

At every batch size, training on $100\%$ correct traces is worse than on $80\%$:
$56.7\to39.5$, $87.1\to59.8$, $87.9\to61.7$, $88.1\to71.7$, $89.2\to79.2$. No
account in this paper predicts this; the drift $\Phi(r)$ is increasing in $\rho$
at every $B$, and so is the supervision-density account of
Section~\ref{sec:competing}.

\paragraph{Run-to-run noise is excluded; a data artefact is not.}
The effect holds in all fifteen cells individually, with per-cell deficits from
$3.4$ to $47.0$ points, so it survives three independent initializations and
batch orderings. It does not follow that the effect is intrinsic: all runs draw
from the same pool, and at $\rho = 1$ that pool is used in its entirety with
correct traces, so the \emph{data} are identical across all fifteen runs and only
initialization and batch order vary. Regenerating the pool under a different
generator seed is the missing control.

\paragraph{A concrete metric confound, which we consider the leading candidate.}
The evaluation described above caps generation at $35$ new tokens. For a word of
length $\ell$ the model must emit a trace of $3\ell-1$ characters, then
\texttt{ : }, then the index: $3\ell+3$ tokens before the answer, plus a newline.
At $\ell = 10$ that is exactly $34$ of the $35$ available, and at $\ell = 9$ it
is $31$. A single spurious token therefore truncates the answer entirely on the
longest words, which are also the words with the lowest Bayes ceiling. Any effect
that makes a model emit longer or less well-calibrated traces will show up as an
accuracy loss under this metric even if the answer-position representation is
correct. At $\rho = 1$ the model has never seen an unpredictable trace and is
maximally committed to generating a full one, which is precisely the regime where
the budget binds hardest. This is testable: under teacher-forced answer scoring,
or with a larger token budget, the $\rho = 1$ deficit should shrink or vanish.

\paragraph{Two substantive candidates, if the confound is ruled out.}
First, \emph{regularisation by corrupted traces}: a minority of corrupted traces
may prevent over-reliance on a brittle copy circuit, forcing the model to
maintain the prompt-only pathway that actually achieves Bayes accuracy. Second,
\emph{optimization pathology at the boundary}: at $\rho = 1$ the trace is
perfectly predictable, trace-token loss falls quickly, and the effective gradient
on the answer token may be reshaped. The answer-only ablation of
Section~\ref{sec:competing} discriminates the second.

We note that the deficits, $17.1, 27.3, 26.2, 16.4, 10.0$, are non-monotone in
$B$, peaking at $B = 64$--$128$. An earlier single-seed version of this table
showed a monotone decay; that trend did not survive replication.

% ==================================================================
\section{Reading the exponent gap}
\label{sec:gap}

We now collect what the theory says about the $-0.72$ measurement. The summary is
that the two most natural objections to the $-1/2$ prediction do not account for
it, that two further mechanisms have the wrong sign, and that what remains is
total sample exposure together with the possibility that three seeds are too few.

\paragraph{Optimizer.}
Section~\ref{sec:adam} carries the analysis through for AdamW under
Assumption~\ref{ass:qs} and obtains $-1/2$, with a numerical-floor correction
pointing the wrong way and of relative size $4\times10^{-6}$
(Remark~\ref{rem:epsnumeric}), and
Theorem~\ref{thm:noexp} extends this to scale-invariant rules at fixed horizon
under transversality. Momentum is a $B$-free constant factor. This does not
account for the gap.

\paragraph{Dimension.}
Corollary~\ref{cor:dimfree} obtains $\Theta(B^{-1/2})$ uniformly in $D$ in a
linear-attention surrogate. This does not account for the gap either, subject to
the surrogate's scope.

\paragraph{Endogenous bias and horizon: wrong sign.}
Proposition~\ref{prop:biasflat} shows bias feedback contributes a
$B$-independent additive term to $\rho_c$, moving the slope toward $0$. Under
AdamW the horizon enters only through the $B$-free $\theta$ of
\eqref{eq:thetadef}.

\paragraph{Sample exposure: the channel we cannot rule out.}
At fixed steps, increasing $B$ increases both the signal-to-noise ratio and the
total number of examples seen, $N = BT$. This is a channel through which $B$
enters outside the minibatch-gradient law, and it is therefore not covered by
Theorem~\ref{thm:noexp} or Remark~\ref{rem:hardfloor}.

\begin{proposition}[Sample-budget susceptibility]
\label{prop:susceptibility}
Suppose the measured threshold is the level set of a criterion of product form
$\Phi(r)\Lambda(N)\ge\theta$, where $\Lambda$ is increasing and represents
progress on a $\rho$-independent component of the task that gates the measured
accuracy---learning the output format, the prompt-search pathway, or reliable
trace generation. Then at fixed $T$,
\begin{equation}
\label{eq:susceptibility}
\frac{d\log\rho_c}{d\log B} = -\frac12 - \chi,
\qquad
\chi := \frac{d\log\Lambda/d\log N}{d\log\Phi/d\log r}\ \ge 0 .
\end{equation}
In the small-$r$ regime $d\log\Phi/d\log r\to1$, so matching $-0.72$ requires
$d\log\Lambda/d\log N\approx0.22$.
\end{proposition}

\begin{proof}
Take logarithms at equality and differentiate with respect to $\log B$ at fixed
$T$, using $\log N = \log B+\log T$ and
$\log r = \log\rho+\tfrac12\log B+\mathrm{const}$.
\end{proof}

A sample-budget exponent of $0.22$ is within the range reported for data-scaling
laws, which we offer as a plausibility check and not as evidence. Propagating the
fifteen-point standard error gives $\chi = 0.22\pm0.05$; propagating the
seed-level error gives $\chi = 0.22\pm0.07$, an interval containing $0$. The
honest position is that this channel is available and quantified, not that it is
established.

\paragraph{Why the obvious control does not work as stated.}
It is tempting to hold $N = BT$ fixed and vary $B$, on the grounds that this
removes the sample-exposure channel. It does not isolate gradient noise. That
protocol sets $T = N/B$, so the horizon changes with $B$, violating (H1) of
Theorem~\ref{thm:noexp} and reintroducing exactly the budget confound of
Remark~\ref{rem:twothresholds}; the cumulative step size $S_T$ changes with it.
Fixed-step, fixed-sample and fixed-epoch protocols identify different compute
allocations, and none isolates minibatch-noise scaling on its own.

\begin{prediction}[Data-diversity control]
\label{pred:diversity}
Fix $B = 256$ and $T = 6000$, and vary only the number of \emph{distinct}
examples in the training pool over $\{N/16,\dots,N\}$, sampling with repetition
so that batch size, step count and cumulative step size are all unchanged. Then
$r$ is unchanged by construction and \eqref{eq:susceptibility} reduces to
$d\log\rho_c/d\log N_{\mathrm{pool}} = -\chi$, so this sweep measures $\chi$
directly. Under the sample-exposure account it should return $\chi\approx0.22$;
under the null that the theory is right and the measured exponent reflects seed
noise, it should return $\chi\approx0$.
\end{prediction}

We prefer this to the fixed-sample-budget sweep because it moves one quantity. A
residual caveat is that repetition changes the effective data distribution, so
the control is clean only if the pool remains large enough to avoid memorisation
at the smallest setting.

% ==================================================================
\section{A competing account, and the experiment that separates it}
\label{sec:competing}

There is a second explanation for a $\rho$-dependent transition with nothing to
do with trace-copy drift, and the standard training convention makes it
available. Under \eqref{eq:lfull} the loss is averaged over both trace and answer
tokens. In our task $|R|\approx3\ell$ and $|Y|\approx1$, so roughly $95\%$ of
supervised tokens are trace tokens. When traces are correct, modelling them
shares structure with answering, so the effort partly transfers; when they are
corrupted they are unpredictable by construction, so the optimizer spends most of
its budget fitting noise and the answer-token signal is diluted. Fitting
incompressible noise on most supervised positions may also degrade the
representations used at the answer position. Under this account the transition
occurs where the useful fraction of the gradient budget crosses a threshold, with
no reference to a trace-copy circuit. Note this account also predicts a
batch-size dependence, so the existence of a batch-size effect does not
discriminate.

The two accounts differ in one respect that can be tested directly.

\begin{definition}[Supervision conditions]
\label{def:conditions}
\textbf{A, full supervision}: $\cL_{\mathrm{full}}$ of \eqref{eq:lfull}.
\textbf{B, answer-only}: the model ingests prompt \emph{and} trace, but
cross-entropy is masked to the answer tokens,
$\cL_{\mathrm{ans}} = |Y|^{-1}\sum_{j\in Y}\cL_{\mathrm{CE}}(y_j,y_j^\star)$.
\textbf{C, direct answer}: no trace in the sequence at all, $W\to Y$.
\end{definition}

Conditions A and B present identical token sequences and differ only in the loss
mask; in our implementation the supervised-token fraction is $41.2\%$ under A and
$4.8\%$ under B. Condition C has no $\rho$ axis: once the trace is removed, a
correct-trace and a corrupted-trace example are byte-identical, so C is a single
baseline number per configuration and should be plotted as a horizontal line.

Under the trace-copy account the mechanism is drift driven by the answer-token
gradient \eqref{eq:gradform}, present in both A and B, so the transition should
persist under B at comparable $\rho_c$. Under the supervision-density account the
mechanism is dilution by trace-token loss, which B removes entirely, so the
transition should shift substantially or vanish. We adopt a shift tolerance of
$0.10$ in $\rho_c$.

Both accounts are testable only if $\rho_c$ is estimable, so we refuse to report
a threshold when the curve rises by less than $15\%$ of the chance-to-ceiling
span, or when the fitted $\rho_c$ falls outside $[-0.25,1.25]$. This guard
matters: an undertrained run produces a flat curve from which a logistic fit will
return a confident but meaningless $\rho_c$.

Condition B has a second benefit given Section~\ref{sec:anomaly}: because the
trace is supplied rather than generated, evaluating under B with the trace in
context removes the generation-length confound as well.

% ==================================================================
\section{Limitations}
\label{sec:limitations}

We list these in descending order of severity.

\paragraph{The evaluation metric is end-to-end, not answer-position.}
Accuracy is measured by free generation from the bare prompt followed by regex
parsing, under a $35$-token cap that binds at $\ell = 10$. This confounds trace
generation, delimiter emission, answer prediction and parsing. It affects the
$\rho = 1$ anomaly directly and, in principle, every number in
Table~\ref{tab:batchsweep}. The primary fix is to report teacher-forced accuracy
at the answer position with the clean trace supplied in context, which is the
quantity the theory is about; we have not run it. That the plateau nonetheless
lands on the computed Bayes ceiling bounds how much the confound can be
distorting the high-$\rho$ regime, but says nothing about the low-$\rho$ or
$\rho = 1$ regimes.

\paragraph{Three seeds are enough to see the effect, not to bound it.}
Three seeds resolved two questions: the $\rho = 1$ anomaly is not run-to-run
noise, and seed variance does peak inside the transition as the theory predicts.
A third remains open: the fifteen-point interval on the exponent excludes $-1/2$
and the seed-level interval does not. Distinguishing them needs more seeds. All
seeds share one data pool, so pool-specific artefacts remain unexcluded
throughout.

\paragraph{The AdamW analysis is a quasi-stationary surrogate.}
Assumption~\ref{ass:qs} requires locally constant gradient statistics over a
window long compared with $(1-\beta_2)^{-1}\approx10^3$ steps, plus a bounded
fourth moment and negligible accumulator covariance. Neither the stationarity nor
the covariance condition holds exactly during warmup or across the transition. A
full treatment would track the joint law of $(u_t,m_t,\nu_t)$. We note that
Theorem~\ref{thm:noexp} does not rely on quasi-stationarity---only on degree-zero
homogeneity, a fixed horizon, uniform third moments and transversality---so the
invariance conclusion survives even if the surrogate constants do not.

\paragraph{The matrix analysis is a linear-attention surrogate.}
Definition~\ref{def:linattn} drops the softmax, and
Assumptions~\ref{ass:alignment} and~\ref{ass:matrix} are assumptions rather than
consequences. The claim established is that dimensionality alone does not alter
the leading exponent in that surrogate; it is not that the scalar reduction is
removed for the four-layer softmax transformer we train.

\paragraph{The segment-selection assumptions are substantive.}
Assumption~\ref{ass:selection} is what makes $\Prob(J=I^\star) = 1/K$ on correct
traces and $1/L$ on corrupted ones. The second clause in particular restricts
what a trained model may do: it could select segments using the corrupted
characters or positional structure, in which case $\mu_0$ and the bias both
change. The assumptions are most plausible near initialization, which is where we
use them.

\paragraph{Per-example independence of $Z$ is a scope condition.}
As Section~\ref{sec:variance} shows, a shared correctness indicator within a
minibatch would leave the between-component variance outside the $1/B$ and the
threshold would not vanish. Our prediction applies to i.i.d.\ mixed pools, not to
regimes where correct and corrupted data arrive in homogeneous blocks---arguably
the more realistic setting for distilled pipelines.

\paragraph{No information gap exists in this task by construction.}
Per Remark~\ref{rem:gap}, our conclusions concern a setting where the trace is
provably useless at the Bayes level. Whether the same picture holds when the
trace genuinely carries information is open, and is the variant we would build
next.

% ==================================================================
\section{Conclusion}
\label{sec:conclusion}

We set out to explain a sharp dependence of accuracy on the fraction of correct
think traces. The first result was negative and clarifying: in this task the
trace is conditionally independent of the answer given the prompt, so the
transition cannot be informational. Both the Bayes ceiling and the chance floor
are computable in closed form, and both are matched by experiment without
fitting.

Relocating the phenomenon to the optimizer, we gave a finite-batch analysis in
which the critical trace fraction is set by a competition between a drift
proportional to $\rho$ and noise proportional to $B^{-1/2}$, with the drift
constant fixed by the task's combinatorics up to a coordinate scale that cancels
from the threshold. The prediction is that the critical fraction vanishes with
batch size as $B^{-1/2}$.

Experiment confirms the vanishing---$\rho_c$ falls by a factor of $7.7$ over a
$16$-fold range of $B$---and measures an exponent of $-0.72\pm0.05$, with every
seed steeper than $-1/2$ but a conservative seed-level interval that contains it.
We then asked whether the two natural objections to the prediction explain the
difference. Under a quasi-stationary surrogate, AdamW changes the mechanism
substantially---a saturating drift, $B$-independent step noise, a
horizon-independent threshold from weight decay---and leaves the exponent at
$-1/2$; in a linear-attention surrogate, parameter dimension leaves it there too.
Neither accounts for the gap, and endogenous bias and horizon coupling have the
wrong sign. What remains is total sample exposure, for which we give a
susceptibility identity and a single-variable control, and the possibility that
three seeds are too few to have measured a real discrepancy at all.

The honest summary is: the qualitative mechanism is probably right; the rate is
probably wrong, though our replication falls short of establishing that; the
candidate explanations are narrowed and quantified rather than settled; and the
measurements that would settle them---more seeds, teacher-forced scoring, and the
diversity control---are specified and runnable.

% ==================================================================
\bibliographystyle{plainnat}
\bibliography{references}

\appendix

% ==================================================================
\section{Proof of Theorem~\ref{thm:escape}}
\label{app:escape}

We prove (i); (ii) is symmetric with inequalities reversed and $\mu_1$ for
$\mu_0$. Let $\cF_t$ be generated by the first $t$ minibatches. On the event that
$u_s\in\cU$ for all $s\le t$, decompose
\[
u_T = u_0 + \underbrace{\sum_{t<T}\eta_t\E[h_t\mid\cF_t]}_{D_T}
+ \underbrace{\sum_{t<T}\eta_t(h_t-\E[h_t\mid\cF_t])}_{M_T}.
\]
In the zero-bias regime Assumption~\ref{ass:main}(1)--(2) give
$\E[h_t\mid\cF_t]\ge\rho\mu_0$, so $D_T\ge\rho\mu_0S_T$. $M_T$ is a martingale
whose predictable quadratic variation is bounded, by \eqref{eq:batchvar}, as
$\langle M\rangle_T = \sum_{t<T}\eta_t^2\Var(h_t\mid\cF_t)\le C_\eta V/B$, and by
Assumption~\ref{ass:main}(4) its increments are sub-exponential with parameter
$K_e$. Freedman's inequality gives
$\Prob(M_T\le-\lambda)\le\exp(-\lambda^2/[2(\nu+K_e\lambda/3)])$ for
$\nu\ge\langle M\rangle_T$.

Since $u_0 = 0$ and escape means $u_T\ge a$, it suffices that $D_T+M_T\ge a$, and
given $D_T\ge\rho\mu_0S_T$ this holds whenever
$M_T\ge a-\rho\mu_0S_T = -\lambda$ with $\lambda := \rho\mu_0S_T-a$, positive
exactly when $\rho>\rho^{+}_{\mathrm{esc}}$. Here $\lambda$ is the actual slack
between cumulative drift and escape radius, not a free parameter; applying
Freedman with $\nu = C_\eta V/B$ gives the stated bound. Since
$\lambda = \mu_0S_T(\rho-\rho^{+}_{\mathrm{esc}})$, separating the two branches
of the denominator gives failure probability at most
\[
\exp\Big(-c\min\Big\{\tfrac{B\mu_0^2S_T^2(\rho-\rho^{+}_{\mathrm{esc}})^2}{C_\eta V},\
\tfrac{\mu_0S_T(\rho-\rho^{+}_{\mathrm{esc}})}{K_e}\Big\}\Big)
\]
for a universal $c$. \qed

% ==================================================================
\section{Proof of Theorem~\ref{thm:matrixescape}}
\label{app:matrix}

\paragraph{Part (i).}
Let $H_t := -\bar G_t$, so $U_{t+1} = U_t+\eta_tH_t$, and let
$u_t := \langle U_t-U_0,\Om^\star\rangle$. The projection is linear and
$\|\Om^\star\|_F = 1$, so $u_{t+1}-u_t = \eta_t\langle H_t,\Om^\star\rangle
=: \eta_th_t$ and $(u_t)$ obeys the scalar recursion. It remains to verify
Assumption~\ref{ass:main} for $h_t$ with $v_+\mapsto v_\parallel$. For the drift,
Assumption~\ref{ass:matrix}(1)--(2) in the zero-bias regime give
$\E[h_t\mid\cF_t] = \rho\mu(U_t)\ge\rho\mu_0$. For the variance, for any unit
$\Om$, $\Var(\langle\hat G,\Om\rangle) =
\langle\Om,\mathrm{Cov}(\vecop\hat G)\Om\rangle\le\langle\Om,\Sigma\Om\rangle$,
so at $\Om^\star$ the per-example projected variance is at most $v_\parallel$;
applying \eqref{eq:totalvar}--\eqref{eq:batchvar} to the projected per-example
gradients---legitimate because they are i.i.d.\ draws from the projected mixture
law by Definition~\ref{def:mixture}---gives
$\Var(h_t\mid\cF_t)\le V_\parallel/B$. Tails follow from
Assumption~\ref{ass:matrix}(4) at $\Om^\star$. Escape from $\cU_a$ in the trace
direction is implied by $u_T\ge a$, since
$\|U_T-U_0\|_F\ge\langle U_T-U_0,\Om^\star\rangle$ by Cauchy--Schwarz, so
Theorem~\ref{thm:escape}(i) applies with $V\mapsto V_\parallel$. No union bound
over directions is taken, hence no dependence on $D$.

\paragraph{Part (ii).}
Confinement requires $\|U_t-U_0\|_F<a$ for all $t\le T$ simultaneously.
Decompose $U_T-U_0 = \mathbf D_T+\mathbf M_T$ with
$\mathbf D_T := \sum_{t<T}\eta_t\E[H_t\mid\cF_t]$ and
$\mathbf M_T := \sum_{t<T}\eta_t(H_t-\E[H_t\mid\cF_t])$, both in $\R^D$. By
Assumption~\ref{ass:matrix}(1), $\|\mathbf D_T\|_F\le\rho\mu_1S_T$.

To apply a matrix martingale inequality to a vector-valued martingale, embed each
increment $\xi_t := \eta_t(H_t-\E[H_t\mid\cF_t])\in\R^D$ in the self-adjoint
dilation
\[
\mathcal D(\xi_t) := \begin{pmatrix}0 & \xi_t^\top\\ \xi_t & 0\end{pmatrix}
\in\R^{(D+1)\times(D+1)},
\]
which satisfies $\|\mathcal D(\xi)\|_{\mathrm{op}} = \|\xi\|_2$ and is linear in
$\xi$, so $\mathcal D(\mathbf M_T) = \sum_t\mathcal D(\xi_t)$ is a self-adjoint
matrix martingale with
$\|\mathcal D(\mathbf M_T)\|_{\mathrm{op}} = \|\mathbf M_T\|_2$. Its predictable
quadratic variation satisfies
\[
\Big\|\sum_{t<T}\E\big[\mathcal D(\xi_t)^2\mid\cF_t\big]\Big\|_{\mathrm{op}}
\;\le\; \sum_{t<T}\eta_t^2\,\Tr\,\mathrm{Cov}(H_t\mid\cF_t)
\;\le\; \frac{C_\eta V_\perp}{B},
\]
using $\mathcal D(\xi)^2 = \mathrm{diag}(\|\xi\|_2^2,\ \xi\xi^\top)$ and bounding
the operator norm of both blocks by the trace. With increments sub-exponential
with parameter $K_e$ in every direction by Assumption~\ref{ass:matrix}(4), the
matrix Freedman inequality of \citet{tropp2011freedman} gives
\[
\Prob\Big(\max_{t\le T}\|\mathbf M_t\|_2\ge\lambda'\Big)
\le (D+1)\exp\Big(-\frac{\lambda'^2}{2(C_\eta V_\perp/B+K_e\lambda'/3)}\Big),
\]
the ambient dimension of the dilation supplying the prefactor. On the complement,
$\|U_t-U_0\|_F\le\|\mathbf D_t\|+\|\mathbf M_t\|<a$ for all $t\le T$ with
$\lambda' = a-\rho\mu_1S_T$. \qed

\paragraph{Where the dimension goes.}
The asymmetry is structural: escaping needs one direction to move far, while
staying needs every direction to stay close. Hence $v_\parallel$ appears in (i)
and $\Tr\Sigma$ in (ii), and the undecided window widens by
$\sqrt{D_{\mathrm{eff}}}$ in high dimension while its endpoints keep the same
batch-size scaling. We have not optimised the $D+1$ factor; a chaining argument
would likely replace it by $\sqrt{\log D_{\mathrm{eff}}}$.

% ==================================================================
\section{Numerical values and calibration}
\label{app:numbers}

For $q = 26$ and $L\sim\mathrm{Unif}\{4,\dots,10\}$, writing $c$ for the gradient
scale of \eqref{eq:gradform}:
\begin{center}
\begin{tabular}{lll}
\toprule
Quantity & Value & Source \\
\midrule
$\bar A^\star$ & $0.8926$ & Proposition~\ref{prop:bayesacc}, averaged over $L$ \\
$A_{\mathrm{chance}} = \E[1/L]$ & $0.1565$ & Corollary~\ref{cor:levels} \\
$\mu_0/c$ & $0.7361$ & Proposition~\ref{prop:mu0} \\
$v_+/c^2 = \E[\tfrac1L(1-\tfrac1L)]$ & $0.1296$ & Lemma~\ref{lem:variance}, $\rho\to0$ \\
$\sqrt{v_+}/\mu_0$ & $0.4892$ & scale-free by Remark~\ref{rem:scale} \\
\bottomrule
\end{tabular}
\end{center}

\paragraph{Calibration.}
Theorem~\ref{thm:scaling} gives $\rho_c(B)\approx0.489\,\kappa/\sqrt B$.
Calibrating at $B = 256$, where the three-seed $\hat\rho_c = 0.175$, gives
$\kappa = 0.175\sqrt{256}/0.489 = 5.72$, after which the theory predicts $\rho_c$
at every other $B$ with no free parameter:
\begin{center}
\begin{tabular}{lccccc}
\toprule
$B$ & 32 & 64 & 128 & 256 & 512 \\
\midrule
predicted $\rho_c$ ($\kappa = 5.72$) & 0.495 & 0.350 & 0.247 & 0.175 & 0.124 \\
measured $\hat\rho_c$ (3 seeds) & 0.759 & 0.465 & 0.305 & 0.175 & 0.099 \\
ratio & 1.53 & 1.33 & 1.23 & 1.00 & 0.80 \\
\bottomrule
\end{tabular}
\end{center}
The monotone trend in the ratio is the exponent discrepancy of
Section~\ref{sec:exponent} in another coordinate; a correct exponent would give a
flat row.

\paragraph{Reading the calibration through the AdamW surrogate.}
Comparing Theorem~\ref{thm:scaling} with \eqref{eq:adamrhoc} at $e_B\approx0$
identifies $\kappa$ with the critical standardized SNR,
$\kappa = \theta/\sqrt{1-\theta^2}$, so $\theta = \kappa/\sqrt{1+\kappa^2} =
0.985$ and $\sqrt{1-\theta^2} = 0.172$. For our schedule the binding branch of
\eqref{eq:thetadef} is the displacement one:
$S_T\approx600\cdot\tfrac12(3{\times}10^{-4}) + 5400\cdot\overline\eta_{\cos}
\approx0.09+0.84 = 0.93$, so $S_T^{-1} = 1.08\gg\lambda_\wdc = 0.01$ and
$\theta = a/S_T$. The calibrated constant therefore corresponds to an escape
radius $a = \theta S_T\approx0.92$ in logit units---the scale on which a softmax
over $L\le10$ positions changes its argmax. We offer this as a plausibility check
and not more: $\theta$ is close to $1$, where $\theta/\sqrt{1-\theta^2}$ is steep
and the inferred $a$ correspondingly sensitive, and the per-seed $\hat\rho_c$ at
$B = 256$ span $0.132$--$0.208$, giving $a$ a range of about $0.89$--$0.94$.

\paragraph{Numerical floor and susceptibility.}
With $\varepsilon_{\mathrm A} = 10^{-8}$ and $\sqrt{v_+}/c = 0.360$,
$e_B = 6.3\times10^{-7}$ at $B = 512$ and $1.6\times10^{-7}$ at $B = 32$; with
$\theta = 0.985$, Theorem~\ref{thm:adamthreshold} predicts a slope of
$-0.5(1-3.7\times10^{-6})$. By Proposition~\ref{prop:susceptibility}, matching
$-0.72$ needs $\chi = 0.22$; over our sweep $N$ ranges from $1.92\times10^5$ to
$3.07\times10^6$, so the required improvement in the $\rho$-independent factor is
$16^{0.22} = 1.8\times$.

% ==================================================================
\section{Deferred material}
\label{app:deferred}

\subsection{Representability of the trace-copy circuit}
\label{app:repr}

The scalar analysis presupposes that the architecture contains a parameterization
implementing the trace-copy family. We keep this separate from the optimization
question and adopt both directions from the literature rather than re-deriving
them.

Standard constructions for two-layer attention-based copying
\citep{olsson2022incontext,ekbote2025what} give sufficiency: if the mechanism
decomposes as one attention operation moving a feature encoding a matched trace
segment into the residual stream at the answer position, and one readout
converting it to answer logits, with positional discrimination required among at
most $2\ell_{\max}+C_0$ positions, then for $m\ge2$ and width
$d\ge C_1\log(2\ell_{\max}+C_0)$ the class contains such a parameterization with
answer error $O(1/T)$ on correctly-traced inputs.

For necessity, \citet{ekbote2025what} show a single attention layer cannot
represent an exact induction-head mechanism for general-order Markov structure
while two layers can, and our task requires the same two-hop composition: locate
a segment by content, then read an index from it. We flag this as an adopted
result whose transfer we have not independently verified---their setting
(arbitrary-order Markov chains) differs from the word-index task in ways that
would need checking---and treat it as an open technical gap.

Together these indicate that $m\ge2$ suffices with adequate width, so the
four-layer architecture is unlikely to face a representability bottleneck. We
stop short of concluding that the transition must therefore be the scalar drift
dynamics: ruling out one alternative does not establish another. The mechanism
could be drift on a trace-copy direction, the supervision-density effect of
Section~\ref{sec:competing}, a curriculum-like ordering effect from a reused
pool, or a combination.

\subsection{Sample complexity of the gradient sign}
\label{app:samplecomplexity}

A complementary measure of difficulty at small $\rho$ is how many i.i.d.\
gradient samples are needed to determine the \emph{sign} of the population drift.

\begin{theorem}
\label{thm:samplecomplexity}
Let Assumption~\ref{ass:main} hold with net drift
$m(\rho) := \rho\mu_0-(1-\rho)\bar\epsilon>0$ at $u_0$. For the empirical mean of
$n$ i.i.d.\ per-example gradients to have the correct sign with probability
$1-\delta$, it suffices that
$n\ge C\max\{\sigma^2\log(1/\delta)/m(\rho)^2,\ K_e\log(1/\delta)/m(\rho)\}$, the
two branches being the Gaussian and sub-exponential regimes of Bernstein's
inequality under Assumption~\ref{ass:main}(4).
\end{theorem}

In the zero-bias regime this reads $n\gtrsim v_+\log(1/\delta)/(\rho^2\mu_0^2)$,
diverging as $\rho^{-2}$: the same qualitative picture as the
information-exponent literature \citep{benarous2021online,abbe2023sgd}, where
weak population signal forces a diverging sample requirement. It is also the same
$\rho^{-2}$ that appears in \eqref{eq:quadratic}: setting $n = B$ recovers
$\rho\gtrsim\sqrt{v_+}/(\mu_0\sqrt B)$. The two results are two views of one
inequality.

\subsection{A reduced mean-field model}
\label{app:meanfield}

The scalar analysis treats one coordinate. The empirical transition involves at
least three interacting quantities, and we sketch a reduced model with the
minimum number of order parameters that can produce a genuine bifurcation. This
is a modelling proposal, not a consequence of the transformer dynamics, which is
why it is deferred here.

Let $x_t\in[0,1]$ be the trace-copy strength (probability the answer readout
attends to a trace segment rather than the prompt), $y_t\in[0,1]$ the
prompt-search strength, and $z_t\ge0$ the match precision, so that the
probability of attending to a segment whose character field equals $C$ is
$\phi(z_t)$ with $\phi(0) = 1/L$ and $\phi(\infty) = 1$. The third parameter is
not redundant: a model can attend to the trace ($x$ large) while matching poorly
($z$ small), in which case it copies an arbitrary index and performs at chance.
It is the product structure that produces accuracy. We model
\begin{align*}
\dot x &= \alpha x(\rho\phi(z)\bar A^\star-\bar p)(1-x)-\alpha\lambda_xx,\\
\dot y &= \beta y(\bar A^\star-\bar p)(1-y)-\beta\lambda_yy,\\
\dot z &= \gamma\rho x(1-\phi(z))-\gamma\lambda_zz,
\end{align*}
with $\bar p$ the mean confidence in the chosen index. Match precision improves
only in proportion to $\rho x$---only when the model both attends to the trace
and the trace is correct---while $x$ grows only when $\phi(z)$ is already large.
This mutual dependence is how the system can have both a stable ``ignore the
trace'' and a stable ``use the trace'' state.

Writing $F_\rho$ for the right-hand side and $(x^0,y^0,z^0)$ for the
trace-ignoring fixed point with $x^0 = z^0 = 0$, call the system non-degenerate
at $\rho^\star$ if $D_{(x,z)}F_{\rho^\star}$ has a simple zero eigenvalue with
right and left eigenvectors $v,w$; $w^\top\partial_\rho F_{\rho^\star}\ne0$; and
$w^\top D^2F_{\rho^\star}[v,v]\ne0$. These are the standard transversality
conditions for a transcritical bifurcation, stated as checkable primitives rather
than assumed via a normal form---the alternative, writing
$\dot x = a(\rho-\rho_c)x-bx^2$ and declaring a bifurcation, assumes the
conclusion. Under them the bifurcation is transcritical and
$x^\star\propto(\rho-\rho^\star)$; if the third condition fails while the cubic
term does not vanish it is a pitchfork with exponent $1/2$; if the second fails
the transition is imperfect and $x^\star$ is smooth. Measuring the exponent of
the normalised accuracy
$\widehat A(\rho) :=
(A(\rho)-A_{\mathrm{chance}})/(\bar A^\star-A_{\mathrm{chance}})$
against $(\rho-\rho_c)$ therefore discriminates the three cases. This requires
denser sampling of $\rho$ near threshold than our grid provides.

\subsection{Beyond word-indexing: a task family}
\label{app:family}

A natural worry is that the transition is a property of word-indexing rather than
of trace supervision. We specify a family of four algorithmic tasks sharing the
surface form $\langle\text{prompt}\rangle\ \langle\text{trace}\rangle\
\texttt{:}\ \langle\text{answer}\rangle$, so the mixing, masking and evaluation
machinery is identical across them.

\begin{center}
\small
\begin{tabular}{llcc}
\toprule
Task & Correct trace & Chance & Ceiling \\
\midrule
\texttt{word\_index} & \texttt{0m 1y 2n \dots} & 0.157 & 0.893 \\
\texttt{count\_char} & running tally \texttt{b0 a1 a2 c2 \dots} & 0.231 & 1.000 \\
\texttt{sort\_letters} & selection sort \texttt{ifkhb>b ifkh>f \dots} & 0.000 & 1.000 \\
\texttt{multiply} & \texttt{4*49=196 60*49=2940 196+2940=3136} & 0.002 & 1.000 \\
\bottomrule
\end{tabular}
\end{center}

Three design points matter. The chance floor is the best constant predictor, not
a uniform guess; for word-index these coincide, since
$\Prob(I^\star = 0) = \E[1/L]$, but for the others they do not, and the constant
predictor is correct because it is what a model ignoring its input achieves.
Getting this wrong inflates the apparent size of a transition. Only word-index
has a ceiling below one, which makes it the informative case for the
reference-level comparison and the others informative for the shape of the
transition. And because the floors differ by two orders of magnitude, cross-task
comparison must use $\widehat A(\rho)$, which places every task on a common
$[0,1]$ scale and leaves $\rho_c$ as the only thing to read off.

% ==================================================================
\section{Reproducibility}
\label{app:repro}

Every sweep reads its defaults from a single configuration module and differs
only in the axis it varies on top of the correct-trace ratio: the ratio sweep,
the task, the batch size at fixed steps, and the batch size at fixed total sample
budget. Dataset pools are generated once and cached, keyed by task, supervision
condition, seed, context length and a fingerprint of the sampler's output, so
editing a task definition invalidates its cache rather than silently reusing data
from an earlier definition.

Two correctness guards are worth recording because both caught real errors during
development. Examples exceeding the context length are counted and reported
rather than silently truncated---a truncated example loses its answer token
entirely and would be scored as a failure with no indication why. And the loss
mask produced by the cached representation is checked for exact equality against
an independently constructed reference mask under all three supervision
conditions of Definition~\ref{def:conditions}.

No equivalent guard existed on the generation budget at evaluation time, which is
how the $35$-token cap of Section~\ref{sec:anomaly} went unnoticed. The
corresponding guard is an assertion that the gold continuation fits within the
budget for every held-out example.

\end{document}
